\chapter{Discrete mathematics, graphs, combinatorics}
\label{vol02:ch17}

\epigraph{The art of programming is the art of organising
complexity, of mastering multitude and avoiding its bastard
chaos as effectively as possible.}{Edsger W. Dijkstra,
\emph{Notes on Structured Programming},
\cite[p.~6]{gen:dijkstra1972}.}

\chapmeta{Half-life: foundational. Archetypes: scaling,
optimisation. Exercise target: 40. Project track: analysis with
optional simulation. Hazard class: none. Reader path default:
standard, with sections explicitly tagged.}

\noindent
The objects of this chapter are countable. Sets of arrangements,
graphs with finitely many vertices, truth tables with finitely
many rows, asymptotic growth classes for algorithms running on
inputs of finite size. The reader who finishes the chapter counts
without enumerating, models a real network as a graph that other
engineers can query, reads $O(n \log n)$ as a contract about how a
program scales, and recognises a combinatorial explosion before the
search has been launched. The mathematics here is the language in
which Volumes VII (Information), VIII (Machines), and IX (Systems)
are written. The chapter is the engineer's working introduction.

\section{Counting, permutations, combinations}\pathtag{core}

A counting problem asks how many objects of a stated kind exist.
The objects might be the orderings of a list, the subsets of a
collection, the words of a fixed length over an alphabet, the paths
of a fixed length in a graph, or the assignments of values to a set
of variables. The discipline is to read the problem carefully
enough to identify which of a small library of counting principles
applies.

\subsection{The two basic rules}

Two rules govern almost every elementary count.

\begin{principle}[Sum rule]
If a set $A$ has $\abs{A}$ elements and a set $B$ disjoint from $A$
has $\abs{B}$ elements, then $\abs{A \cup B} = \abs{A} + \abs{B}$.
The rule extends to any finite collection of pairwise disjoint sets.
\end{principle}

\begin{principle}[Product rule]
If a process consists of two stages, the first having $m$ outcomes
and the second, regardless of the first stage's outcome, having $n$
outcomes, then the process has $m \cdot n$ outcomes. The rule
extends to any finite sequence of stages.
\end{principle}

The product rule licenses the count of $26^{8}$ for the number of
length-eight strings over the Latin alphabet, the count of $10^{6}$
for the number of decimal six-digit numbers (including those with
leading zeros), the count of $2^{n}$ for the number of subsets of
an $n$-element set, and the count of $26 \cdot 10^{3}$ for the
number of automobile licence plates of the form letter-digit-digit-digit
in a jurisdiction that uses that template. Each comes from a stage
count: stage one has $26$ outcomes, stage two has $10$, stage three
has $10$, stage four has $10$, multiply.

The two rules combine. A six-character password drawn from twenty-six
lowercase letters or ten digits has $36^{6}$ choices total. The
count of length-six strings that contain at least one digit is more
work: count all length-six strings over the combined alphabet,
subtract the length-six strings using letters only.
\[
36^{6} - 26^{6} = 2{,}176{,}782{,}336 - 308{,}915{,}776 =
1{,}867{,}866{,}560.
\]
The complement principle, count the unwanted and subtract from the
total, is the third routine tool of the elementary count.

\subsection{Permutations}

A \engterm{permutation} of an $n$-element set is an ordered
arrangement of its elements. The number of permutations of $n$
distinct objects is
\[
n! = n \cdot (n-1) \cdot (n-2) \cdots 2 \cdot 1,
\]
read $n$-factorial. The product rule supplies the proof: there are
$n$ choices for the first position, $n-1$ choices for the second
once the first is taken, and so on down to one choice for the last.

A \engterm{$k$-permutation} of an $n$-element set is an ordered
arrangement of $k$ of its elements. The count is
\[
P(n, k) = \frac{n!}{(n-k)!} = n \cdot (n-1) \cdots (n-k+1).
\]
The product is read $n$ choices for position one, $n-1$ for position
two, and so on for $k$ positions. The factorial form is the same
quantity rewritten. Engineers carry both because the factorial
ratio cancels in symbolic work and the product form computes
without ever forming the large numerator.

A common variant counts permutations of objects with repetition. If
$n$ objects fall into $k$ classes of sizes $n_{1}, n_{2}, \ldots,
n_{k}$ with $n_{1} + \cdots + n_{k} = n$, and objects in the same
class are indistinguishable, the count of distinct arrangements is
the \engterm{multinomial coefficient}
\[
\binom{n}{n_{1}, n_{2}, \ldots, n_{k}} =
\frac{n!}{n_{1}!\, n_{2}! \cdots n_{k}!}.
\]
The number of distinct arrangements of the letters of MISSISSIPPI
(four S, four I, two P, one M) is therefore $11! / (4! \cdot 4!
\cdot 2! \cdot 1!) = 34{,}650$.

\subsection{Combinations}

A \engterm{combination} of $k$ objects from $n$ is an unordered
subset of size $k$. The count is the binomial coefficient
\[
\binom{n}{k} = \frac{n!}{k!\,(n-k)!},
\]
which is the number of $k$-permutations divided by the $k!$
orderings of any one combination.

The binomial coefficient counts the five-card poker hands as
$\binom{52}{5} = 2{,}598{,}960$, the bridge hands as $\binom{52}{13}
= 635{,}013{,}559{,}600$, the lottery tickets in a six-of-forty-nine
draw as $\binom{49}{6} = 13{,}983{,}816$, and the binary strings of
length $n$ with exactly $k$ ones as $\binom{n}{k}$. The reader who
internalises the binomial coefficient as ``the count of $k$-out-of-$n$
choices'' will read the next four chapters of this volume without
re-deriving the form.

Three identities recur often enough that an engineer holds them
without lookup.

\begin{principle}[Pascal's identity]
\[
\binom{n}{k} = \binom{n-1}{k-1} + \binom{n-1}{k}.
\]
A combinatorial proof: a $k$-subset of $\set{1, \ldots, n}$ either
contains element $n$ (the rest is a $(k-1)$-subset of $\set{1,
\ldots, n-1}$) or does not (the rest is a $k$-subset of $\set{1,
\ldots, n-1}$). The two cases are disjoint and exhaustive; the sum
rule supplies the identity.
\end{principle}

\begin{principle}[Binomial theorem]
For commuting $a, b$ and any non-negative integer $n$,
\[
(a + b)^{n} = \sum_{k=0}^{n} \binom{n}{k}\, a^{n-k}\, b^{k}.
\]
\end{principle}

Setting $a = b = 1$ in the binomial theorem gives $\sum_{k=0}^{n}
\binom{n}{k} = 2^{n}$, the count of all subsets of an $n$-element
set, which is also a direct product-rule count: each element is
independently in or out.

\begin{principle}[Vandermonde's identity]
\[
\binom{m+n}{r} = \sum_{k=0}^{r} \binom{m}{k}\binom{n}{r-k}.
\]
The combinatorial reading: choose $r$ from a set of $m+n$ objects
split into a group of $m$ and a group of $n$, summing over the
number drawn from the first group.
\end{principle}

\subsection{The pigeonhole principle}

The pigeonhole principle is the simplest theorem in combinatorics
that earns its name in real engineering arguments.

\begin{principle}[Pigeonhole]
If $n$ objects are placed into $k$ boxes and $n > k$, then some box
contains at least two objects. More generally, if $n$ objects are
placed into $k$ boxes, some box contains at least $\lceil n/k \rceil$
objects.
\end{principle}

The proof is contradiction: if every box contained at most one
(more generally, at most $\lceil n/k \rceil - 1$) object, the total
would not exceed $k$ (more generally, $k(\lceil n/k \rceil - 1) < n$),
contradicting the count $n$.

The principle solves problems whose statement does not announce it.
Among any thirteen people, two share a birth month. Among any $n+1$
integers chosen from $\set{1, 2, \ldots, 2n}$, two are coprime.
Among any $n+1$ integers in $\set{1, 2, \ldots, 2n}$, one divides
another. In a hash table of $k$ buckets accepting $n$ keys, the most
loaded bucket holds at least $\lceil n/k \rceil$ keys, which is the
worst-case-collision bound an engineer cites when sizing a hash
table.

\begin{archetype}[Scaling]
A counting problem scales with its parameters in a way the engineer
can read off the form. A factorial $n!$ grows faster than any
exponential $c^{n}$ for fixed $c$; an exponential $2^{n}$ grows
faster than any polynomial $n^{k}$ for fixed $k$; a polynomial
$n^{k}$ grows faster than any logarithm $\log^{j} n$ for fixed $j$.
The order, polynomial dominated by exponential dominated by
factorial, is the engineer's first scaling literacy. The reader
will meet the same hierarchy in section 17.5 as the order of growth
of running times, in Chapter 18 as the curse of dimensionality, and
in Volume VIII as the size of state spaces an algorithm cannot
afford to enumerate.
\end{archetype}

\subsection{Inclusion-exclusion}

The sum rule gives $\abs{A \cup B} = \abs{A} + \abs{B}$ when $A$
and $B$ are disjoint. When they overlap, the count $\abs{A} + \abs{B}$
double-counts the intersection, so
\[
\abs{A \cup B} = \abs{A} + \abs{B} - \abs{A \cap B}.
\]

For three sets,
\[
\abs{A \cup B \cup C} = \abs{A} + \abs{B} + \abs{C} - \abs{A \cap B}
- \abs{A \cap C} - \abs{B \cap C} + \abs{A \cap B \cap C}.
\]

The general statement, the \engterm{inclusion-exclusion principle},
is
\[
\Big| \bigcup_{i=1}^{n} A_{i} \Big| = \sum_{k=1}^{n}
(-1)^{k+1} \sum_{1 \leq i_{1} < \cdots < i_{k} \leq n}
\abs{A_{i_{1}} \cap \cdots \cap A_{i_{k}}}.
\]
The alternating sign accounts for the over- and under-counting of
each intersection across the union.

The engineer's working level is recognition rather than facility:
when a problem asks for the count of objects satisfying at least
one of several properties, inclusion-exclusion is the framework
even if the case count is small. A typical engineering instance is
the count of permutations with no fixed points (the number of
derangements), which inclusion-exclusion delivers as $D_{n} = n!
\sum_{k=0}^{n} (-1)^{k}/k!$, a quantity the reader will meet again
in queueing analysis and in error-correcting code design.

\subsection{Estimation before calculation}

\begin{estimation}
\textbf{Question.} A network operator runs a distributed system
across $n = 20$ data centres. Each pair of data centres can be
connected directly by a private link, or routed through one or more
intermediaries. Estimate, before reading on, the number of distinct
fully-meshed connection topologies, that is, the number of ways to
select a non-empty subset of the $\binom{20}{2}$ possible links and
declare those to be the operational topology.

\smallskip
\textit{Estimate, before reading on.}

\smallskip
The number of unordered pairs is $\binom{20}{2} = 190$. Each link
is independently in or out, giving $2^{190}$ topologies, minus the
single empty topology if we require non-emptiness, so $2^{190} - 1$.
$2^{190}$ has roughly $190 \cdot \log_{10} 2 \approx 57$ decimal
digits.

The estimate matters because a search over topologies cannot be
exhaustive: $10^{57}$ is many orders of magnitude beyond the number
of operations a current data centre can run in a human lifetime
(which is roughly $10^{18}$ to $10^{20}$). A network designer who
proposes to ``try every topology'' has misread the problem. Section
17.8 names this miscalibration as the dominant failure mode for
combinatorial design spaces.
\end{estimation}

\section{Graph theory: nodes, edges, walks, paths}\pathtag{core}

A graph is the engineer's working data structure for any system
whose objects relate pairwise. Power grids, communication networks,
transportation networks, social networks, software dependencies,
chemical reaction networks, and the call graph of a running program
are graphs. The vocabulary in this section is the vocabulary the
reader needs to read the rest of the chapter and to read the
algorithms literature.

\subsection{Definitions}

\begin{definition}[Graph]
A \engterm{graph} $G = (V, E)$ consists of a finite set $V$ of
\engterm{vertices} (also called \engterm{nodes}) and a set $E$ of
\engterm{edges}. In an \engterm{undirected} graph each edge is an
unordered pair $\set{u, v}$ with $u \neq v$. In a \engterm{directed}
graph (also called a \engterm{digraph}), each edge is an ordered
pair $(u, v)$, with $u$ the tail and $v$ the head. A
\engterm{weighted} graph carries a function $w: E \to \R$ assigning
a real-valued weight to each edge.
\end{definition}

We write $n = \abs{V}$ for the number of vertices and $m = \abs{E}$
for the number of edges. The two natural extreme regimes are the
\engterm{sparse} regime $m = O(n)$, where each vertex has on
average a constant number of neighbours, and the \engterm{dense}
regime $m = \Theta(n^{2})$, where the edge count grows with the
square of the vertex count. Most engineering networks are sparse:
a power grid bus connects to a small constant number of lines, a
city intersection touches a small constant number of streets.

The \engterm{degree} of a vertex $v$ in an undirected graph is the
number of edges incident to $v$, written $\deg(v)$. In a digraph,
the \engterm{in-degree} $\deg^{-}(v)$ counts edges with $v$ as head
and the \engterm{out-degree} $\deg^{+}(v)$ counts edges with $v$ as
tail.

\begin{principle}[Handshake lemma]
In any undirected graph,
\[
\sum_{v \in V} \deg(v) = 2 \abs{E}.
\]
The proof is a count: each edge contributes to the degree of each
of its two endpoints. The corollary is that the number of
odd-degree vertices in any undirected graph is even.
\end{principle}

\subsection{Walks, paths, cycles}

A \engterm{walk} from $u$ to $v$ is a sequence $u = v_{0}, v_{1},
\ldots, v_{k} = v$ of vertices such that each consecutive pair
$\set{v_{i}, v_{i+1}}$ is an edge. The walk's \engterm{length} is
$k$ (the number of edges traversed). A \engterm{trail} is a walk
with no repeated edges. A \engterm{path} is a walk with no repeated
vertices. A \engterm{cycle} is a walk that begins and ends at the
same vertex, has no repeated edges, and (apart from the shared
endpoint) no repeated vertices.

In a directed graph, all of the above respect edge orientation: a
directed walk from $u$ to $v$ traverses edges with their
orientation, and a directed cycle returns to its starting vertex
along oriented edges. A directed graph that contains no directed
cycle is a \engterm{directed acyclic graph} (DAG); DAGs are the
canonical model for dependency, precedence, and irreversibility,
and the canonical structure underlying topological sort and
critical-path analysis (Volume IX).

\subsection{Connectedness, components, trees}

An undirected graph is \engterm{connected} if there is a path
between every pair of vertices. A \engterm{connected component} of
a graph is a maximal connected subgraph. The vertex set partitions
into connected components.

A directed graph is \engterm{strongly connected} if for every
ordered pair $(u, v)$ there is a directed path from $u$ to $v$. A
\engterm{strongly connected component} is a maximal strongly
connected subgraph; the components form the vertices of a DAG (the
\engterm{condensation}) on which higher-level reasoning takes
place.

\begin{definition}[Tree]
A \engterm{tree} is a connected undirected graph with no cycles.
Equivalently, a tree on $n$ vertices is a connected graph with
exactly $n - 1$ edges. A \engterm{forest} is a disjoint union of
trees.
\end{definition}

Trees admit a root, after which every non-root vertex has a unique
parent and every vertex has a (possibly empty) set of children. The
\engterm{depth} of a vertex is the length of the path from the root
to the vertex; the \engterm{height} of the tree is the maximum
depth. The reader will meet trees again as parse trees in compilers
(Volume VII), as decision trees in machine learning (Volume X), as
binary search trees in algorithms (Volume VIII), and as spanning
trees in section 17.3 of this chapter.

\subsection{Working metrics}

Three numbers describe the gross structure of a graph and recur
across applications.

\begin{definition}[Diameter]
The \engterm{distance} $d(u, v)$ between two vertices in a
connected graph is the length of the shortest path between them
(in a weighted graph, the minimum weight of any path). The
\engterm{diameter} of the graph is $\max_{u, v} d(u, v)$.
\end{definition}

\begin{definition}[Clustering coefficient]
The \engterm{local clustering coefficient} of a vertex $v$ with
$\deg(v) \geq 2$ is the fraction of pairs of $v$'s neighbours that
are themselves adjacent:
\[
C(v) = \frac{2 \abs{\set{(u, w) : u, w \in N(v),\, \set{u, w} \in
E}}}{\deg(v)(\deg(v) - 1)},
\]
where $N(v)$ is the neighbour set of $v$. The \engterm{global
clustering coefficient} is the average of $C(v)$ over the vertex
set, restricted to vertices with $\deg(v) \geq 2$.
\end{definition}

The clustering coefficient measures the density of triangles in the
graph and so the degree to which ``my friend's friend is also my
friend.'' Real social and biological networks routinely have high
clustering relative to a random graph of the same edge count.

The \engterm{degree distribution} $P(k)$ is the probability that a
randomly chosen vertex has degree $k$. Many real networks exhibit a
heavy-tailed degree distribution, often well-approximated by a
power law $P(k) \propto k^{-\gamma}$ over a working range. The
heavy tail is what makes a small set of high-degree \engterm{hubs}
the engineer's first concern when reasoning about robustness:
removing a random vertex barely affects connectivity, removing a
hub fragments the graph.

\subsection{Representations}

A graph $G = (V, E)$ has two standard digital representations.

The \engterm{adjacency matrix} is the $n \times n$ matrix $A$ with
$A_{ij} = 1$ if $\set{i, j} \in E$ (or, in the weighted case,
$A_{ij} = w(i, j)$) and $A_{ij} = 0$ otherwise. The matrix is
symmetric for undirected graphs. Storage is $\Theta(n^{2})$,
independent of $m$.

The \engterm{adjacency list} represents $G$ as an array of length
$n$ in which entry $i$ is the list of neighbours of vertex $i$.
Storage is $\Theta(n + m)$.

For sparse graphs ($m = O(n)$ or $m = O(n \log n)$), the adjacency
list is overwhelmingly the engineer's choice. For dense graphs and
for algorithms that exploit matrix algebra (powers of $A$ count
walks; eigenvalues of $A$ encode connectivity), the adjacency
matrix is the working representation. The choice is rarely a wash;
it is the first design decision in implementing a graph algorithm.

\subsection{Worked example: a small social graph}

Consider a graph $G$ on $V = \set{1, 2, 3, 4, 5}$ with edges $E =
\set{\set{1,2}, \set{1,3}, \set{2,3}, \set{2,4}, \set{3,4},
\set{4,5}}$, drawn in figure
\ref{fig:vol02:ch17:worked-graph}. The graph has $n = 5$ vertices
and $m = 6$ edges. Degrees: $\deg(1) = 2, \deg(2) = 3, \deg(3) = 3,
\deg(4) = 3, \deg(5) = 1$. The handshake lemma checks: the degree
sum is $2 + 3 + 3 + 3 + 1 = 12 = 2 m$.

\begin{figure}[ht]
\centering
\begin{tikzpicture}[
  vertex/.style = {circle, draw=engblack, thick, fill=englime!25,
                   inner sep=0pt, minimum size=7mm, font=\small},
  edge/.style   = {thick, draw=engblack},
  hub/.style    = {circle, draw=engblue, very thick, fill=engblue!15,
                   inner sep=0pt, minimum size=7mm, font=\small},
]
% Worked five-vertex social graph used in section 17.2 and exercise 17.3.
\node[vertex] (v1) at (0, 2)     {$1$};
\node[vertex] (v2) at (1.8, 3)   {$2$};
\node[vertex] (v3) at (1.8, 1)   {$3$};
\node[hub]    (v4) at (3.8, 2)   {$4$};
\node[vertex] (v5) at (5.8, 2)   {$5$};

\draw[edge] (v1) -- (v2);
\draw[edge] (v1) -- (v3);
\draw[edge] (v2) -- (v3);
\draw[edge] (v2) -- (v4);
\draw[edge] (v3) -- (v4);
\draw[edge] (v4) -- (v5);

% Degree annotations
\node[font=\scriptsize, anchor=east] at ($(v1.west) + (-0.1, 0)$) {$\deg=2$};
\node[font=\scriptsize, anchor=south] at ($(v2.north) + (0, 0.1)$) {$\deg=3$};
\node[font=\scriptsize, anchor=north] at ($(v3.south) + (0, -0.1)$) {$\deg=3$};
\node[font=\scriptsize, anchor=south, engblue]
  at ($(v4.north) + (0, 0.15)$) {hub, $\deg=3$};
\node[font=\scriptsize, anchor=west] at ($(v5.east) + (0.1, 0)$) {$\deg=1$, leaf};
\end{tikzpicture}
\caption{The worked five-vertex graph of section 17.2. Vertex $4$
carries the only edge to vertex $5$ and is the graph's articulation
point: its removal disconnects vertex $5$ from the rest. The
clustering coefficient is high at vertices $2$ and $3$ ($C = 2/3$
each) because the triangle $\set{1, 2, 3}$ closes their neighbour
pairs.}
\label{fig:vol02:ch17:worked-graph}
\end{figure}


The graph is connected: every vertex reaches every other along the
edge set. The diameter is $d(1, 5) = d(3, 5) = 2$ (each goes through
vertex $4$); other distances are at most $2$, so the diameter is
$2$. The clustering coefficient at vertex $2$ is the fraction of
edges among its neighbours $\set{1, 3, 4}$: edges $\set{1,3}$ and
$\set{3,4}$ exist, $\set{1,4}$ does not, so $C(2) = 2/3$. The same
calculation at vertex $3$ gives $C(3) = 2/3$. At vertex $4$, the
neighbours are $\set{2, 3, 5}$; the only neighbour-pair edge is
$\set{2, 3}$, giving $C(4) = 1/3$.

Vertex $4$ also illustrates the \engterm{articulation-point}
concept: its removal disconnects vertex $5$ from the rest of the
graph. The articulation points and \engterm{bridges} (edges whose
removal disconnects the graph) of a network are the single-point
failures an engineer locates before commissioning a system. Both
can be found in $O(n + m)$ time by Tarjan's depth-first algorithm
on the graph's DFS tree, the standard pre-flight check on any
network whose service depends on connectivity.

\subsection{Special graph families}\pathtag{standard}

Three families recur often enough that an engineer reads them on
sight (figure \ref{fig:vol02:ch17:graph-families}).

\begin{figure}[ht]
\centering
\begin{tikzpicture}[
  vertex/.style = {circle, draw=engblack, thick, fill=englime!25,
                   inner sep=0pt, minimum size=5mm, font=\scriptsize},
  bvertex/.style = {circle, draw=engblack, thick, fill=engblue!15,
                    inner sep=0pt, minimum size=5mm, font=\scriptsize},
  edge/.style   = {thick, draw=engblack},
]

% --- Panel 1: planar / non-planar (left)
\begin{scope}[xshift=0cm]
\node[font=\small, anchor=south] at (1.5, 3.4) {(a) planar};
\node[vertex] (a1) at (0, 2)    {};
\node[vertex] (a2) at (3, 2)    {};
\node[vertex] (a3) at (1.5, 0)  {};
\node[vertex] (a4) at (1.5, 3)  {};
\draw[edge] (a1) -- (a2);
\draw[edge] (a1) -- (a3);
\draw[edge] (a2) -- (a3);
\draw[edge] (a1) -- (a4);
\draw[edge] (a2) -- (a4);
\end{scope}

% --- Panel 2: bipartite (middle)
\begin{scope}[xshift=5cm]
\node[font=\small, anchor=south] at (1.5, 3.4) {(b) bipartite $K_{3,3}$};
\foreach \i/\name in {0/u1, 1/u2, 2/u3} {
  \node[vertex] (\name) at (\i*1.5, 3) {};
}
\foreach \i/\name in {0/v1, 1/v2, 2/v3} {
  \node[bvertex] (\name) at (\i*1.5, 0) {};
}
\foreach \src in {u1, u2, u3}
  \foreach \dst in {v1, v2, v3}
    \draw[edge] (\src) -- (\dst);
\end{scope}

% --- Panel 3: tree (right)
\begin{scope}[xshift=10cm]
\node[font=\small, anchor=south] at (1.5, 3.4) {(c) tree};
\node[vertex] (r)  at (1.5, 3)   {};
\node[vertex] (c1) at (0.3, 2)   {};
\node[vertex] (c2) at (2.7, 2)   {};
\node[vertex] (g1) at (-0.3, 1)  {};
\node[vertex] (g2) at (0.9, 1)   {};
\node[vertex] (g3) at (2.1, 1)   {};
\node[vertex] (g4) at (3.3, 1)   {};
\node[vertex] (h1) at (-0.3, 0)  {};
\node[vertex] (h2) at (0.9, 0)   {};
\draw[edge] (r) -- (c1);
\draw[edge] (r) -- (c2);
\draw[edge] (c1) -- (g1);
\draw[edge] (c1) -- (g2);
\draw[edge] (c2) -- (g3);
\draw[edge] (c2) -- (g4);
\draw[edge] (g1) -- (h1);
\draw[edge] (g1) -- (h2);
\end{scope}

\end{tikzpicture}
\caption{Three graph families the engineer reads on sight.
\textbf{(a)} A planar graph drawable without edge crossings;
the complete graph $K_{4}$ here is the smallest non-trivial planar
graph that fills its faces. \textbf{(b)} The complete bipartite
$K_{3,3}$, the canonical non-planar graph and a Kuratowski
obstruction; bipartite structure makes it a stock model for
assignment problems. \textbf{(c)} A tree on nine vertices, eight
edges, height three; root at top.}
\label{fig:vol02:ch17:graph-families}
\end{figure}


A \engterm{planar} graph is one that can be drawn in the plane
without edge crossings. Planarity is the structural property a
printed-circuit-board layout, a road map, and a chemical-bond
diagram all enjoy in the limit of no via, no overpass, no
out-of-plane bond. Kuratowski's theorem (1930) gives the test: a
graph is planar if and only if it contains no subgraph that is a
subdivision of $K_{5}$ (the complete graph on five vertices) or
$K_{3,3}$ (the complete bipartite graph on $3 + 3$ vertices).
Euler's formula $v - e + f = 2$ for a connected planar graph (with
$v$ vertices, $e$ edges, $f$ faces including the outer face) gives
the upper bound $e \leq 3v - 6$ on edge count, the engineer's
working ceiling for a planar topology. The four-colour theorem
states that every planar graph admits a proper vertex colouring
with at most four colours; the result is used in frequency
assignment, register allocation, and scheduling problems whose
conflict graphs are planar. We refer the reader to
\cite{text:v2ch17-bondy-murty} for the proofs.

A \engterm{bipartite} graph $G = (V, E)$ has $V = A \cup B$ with
$A \cap B = \emptyset$ and every edge between $A$ and $B$. A graph
is bipartite if and only if it has no odd-length cycle, a test that
runs in $O(n + m)$ time by two-colouring during a BFS. The
bipartite case is the canonical structure for assignment problems
(section 17.3 matching), for tracking which workers can perform
which tasks, which channels are compatible with which receivers,
and which test cases exercise which code modules.

A \engterm{tree} (introduced above) has the redundant
characterisation that it is the unique connected graph with no
cycles, equivalently the connected graph with $n - 1$ edges.
\engterm{Rooted trees} carry a parent-child hierarchy and are the
working data structure for hierarchical organisation: file systems,
parse trees, decision trees, dependency trees, and the
divide-and-conquer recursion trees we meet in section 17.6.

\subsection{Graph traversal: BFS and DFS}\pathtag{standard}

Almost every graph algorithm reduces, in part, to a systematic
exploration of vertices reachable from a start vertex. Two
traversal disciplines, breadth-first and depth-first, are the
engineer's working primitives (figure
\ref{fig:vol02:ch17:bfs-dfs}).

\begin{figure}[ht]
\centering
\begin{tikzpicture}[
  vertex/.style = {circle, draw=engblack, thick, fill=englime!25,
                   inner sep=0pt, minimum size=7mm, font=\small},
  visited/.style = {circle, draw=engblue, very thick, fill=engblue!15,
                    inner sep=0pt, minimum size=7mm, font=\small},
  edge/.style = {thick, draw=enggray},
  tedge/.style = {very thick, draw=engblue},
]

% Shared underlying graph.
% Layout: root A at top, two children B (left), C (right),
% A-B-D-F branch and A-C-E-G branch with cross edges.

% Panel 1: BFS
\begin{scope}[xshift=0cm]
\node[font=\small, anchor=south] at (2, 4.6) {(a) BFS from $A$};
\node[visited] (A1) at (2, 4)   {$A$};
\node[visited] (B1) at (0.5, 3) {$B$};
\node[visited] (C1) at (3.5, 3) {$C$};
\node[visited] (D1) at (-0.5, 2) {$D$};
\node[visited] (E1) at (2, 2)    {$E$};
\node[visited] (F1) at (-0.5, 1) {$F$};
\node[visited] (G1) at (2, 1)    {$G$};

\draw[tedge] (A1) -- (B1);
\draw[tedge] (A1) -- (C1);
\draw[tedge] (B1) -- (D1);
\draw[tedge] (C1) -- (E1);
\draw[tedge] (D1) -- (F1);
\draw[tedge] (E1) -- (G1);
\draw[edge]  (C1) -- (D1);
\draw[edge]  (B1) -- (E1);

% Discovery order along the side.
\node[font=\scriptsize, anchor=west, engblue] at (4.3, 4) {$1$};
\node[font=\scriptsize, anchor=east, engblue] at (-0.1, 3) {$2$};
\node[font=\scriptsize, anchor=west, engblue] at (4.1, 3) {$3$};
\node[font=\scriptsize, anchor=east, engblue] at (-1.1, 2) {$4$};
\node[font=\scriptsize, anchor=south, engblue] at (2, 2.4) {$5$};
\node[font=\scriptsize, anchor=east, engblue] at (-1.1, 1) {$6$};
\node[font=\scriptsize, anchor=south, engblue] at (2, 1.4) {$7$};
\node[font=\footnotesize, anchor=north]
  at (2, 0.2) {queue $\to$ level-order};
\end{scope}

% Panel 2: DFS
\begin{scope}[xshift=8cm]
\node[font=\small, anchor=south] at (2, 4.6) {(b) DFS from $A$};
\node[visited] (A2) at (2, 4)   {$A$};
\node[visited] (B2) at (0.5, 3) {$B$};
\node[visited] (C2) at (3.5, 3) {$C$};
\node[visited] (D2) at (-0.5, 2) {$D$};
\node[visited] (E2) at (2, 2)    {$E$};
\node[visited] (F2) at (-0.5, 1) {$F$};
\node[visited] (G2) at (2, 1)    {$G$};

\draw[tedge] (A2) -- (B2);
\draw[tedge] (B2) -- (D2);
\draw[tedge] (D2) -- (F2);
\draw[tedge] (B2) -- (E2);
\draw[tedge] (E2) -- (G2);
\draw[tedge] (A2) -- (C2);
\draw[edge]  (C2) -- (D2);
\draw[edge]  (C2) -- (E2);

\node[font=\scriptsize, anchor=west, engblue] at (4.3, 4) {$1$};
\node[font=\scriptsize, anchor=east, engblue] at (-0.1, 3) {$2$};
\node[font=\scriptsize, anchor=east, engblue] at (-1.1, 2) {$3$};
\node[font=\scriptsize, anchor=east, engblue] at (-1.1, 1) {$4$};
\node[font=\scriptsize, anchor=south, engblue] at (2, 2.4) {$5$};
\node[font=\scriptsize, anchor=south, engblue] at (2, 1.4) {$6$};
\node[font=\scriptsize, anchor=west, engblue] at (4.1, 3) {$7$};
\node[font=\footnotesize, anchor=north]
  at (2, 0.2) {stack $\to$ depth-first};
\end{scope}

\end{tikzpicture}
\caption{Two traversal orders on the same graph. BFS \textbf{(a)}
visits vertices in increasing distance from the source, the
guarantee that makes BFS the right shortest-path algorithm on
unweighted graphs. DFS \textbf{(b)} drives along one branch as
far as it can before backtracking, the discipline that makes DFS
the right algorithm for topological sort, cycle detection, and
strongly-connected-component decomposition. The choice of queue
(BFS) or stack (DFS) is the single line that distinguishes their
implementations.}
\label{fig:vol02:ch17:bfs-dfs}
\end{figure}


\engterm{Breadth-first search} (BFS) maintains a queue of vertices
to visit, dequeues a vertex, enqueues each of its undiscovered
neighbours, and records the discovery order. On an unweighted
graph, the discovery order is a level order from the source, so
BFS computes the unweighted shortest-path distance from the source
to every reachable vertex in $\Theta(n + m)$ time. BFS is the
right answer for unweighted shortest paths, connectivity testing,
bipartite-graph two-colouring, and any problem that asks ``what is
the smallest number of edges between $u$ and $v$.''

\engterm{Depth-first search} (DFS) replaces the queue with a stack
(or with recursion). The traversal drives along one branch as far
as the edges permit, backtracks when it runs out, and tries the
next branch. DFS runs in $\Theta(n + m)$ and yields the structural
information on which a family of standard algorithms is built:
topological sort of a DAG, cycle detection in directed and
undirected graphs, strongly-connected-component decomposition
(Tarjan and Kosaraju), articulation-point and bridge identification,
and the back-edge / cross-edge / forward-edge classification that
underwrites the analysis of recursive procedures on graphs.

The choice between BFS and DFS is rarely a wash. The engineer
picks BFS when the application is about distance from a source, the
shallowest match, or level structure; DFS when the application is
about reachability, ordering, or recursive structure on the
exploration tree. The implementations differ in one line of code
and produce algorithmically distinct outputs.

\section{Algorithmic graph problems}\pathtag{standard}

A graph by itself does not solve an engineering problem. Algorithms
do. Three algorithmic problems and their working solutions form the
backbone of network engineering: shortest path, minimum spanning
tree, and matching. We present the algorithms in the form an
engineer reads off CLRS \cite{text:cormen-clrs} or
Kleinberg-Tardos \cite{text:kleinberg-tardos}: pseudocode, running
time in big-$O$ form, and the precondition that distinguishes the
algorithm's correctness regime.

\subsection{Shortest path: Dijkstra}

Given a directed weighted graph $G = (V, E)$ with non-negative edge
weights $w(u, v) \geq 0$ and a source vertex $s$, the
\engterm{single-source shortest-path problem} asks for the shortest
distance $d(s, v)$ from $s$ to every other vertex $v$.

\begin{principle}[Dijkstra's algorithm]
\textbf{Input}: digraph $G = (V, E)$, weights $w \geq 0$, source $s$.
\textbf{Output}: distances $d(s, v)$ for every $v \in V$.

\begin{enumerate}[leftmargin=*]
  \item Initialise $d(s, s) = 0$ and $d(s, v) = \infty$ for every
    $v \neq s$. Mark every vertex \emph{unvisited}.
  \item While an unvisited vertex with $d(s, \cdot) < \infty$
    exists: select the unvisited vertex $u$ with the smallest $d(s,
    u)$. For every edge $(u, v) \in E$, if $d(s, u) + w(u, v) <
    d(s, v)$, set $d(s, v) = d(s, u) + w(u, v)$. Mark $u$
    visited.
  \item Return the distance vector.
\end{enumerate}
\end{principle}

The algorithm's correctness rests on the loop invariant that, when
$u$ is selected, $d(s, u)$ is the true shortest-path distance from
$s$ to $u$. The invariant fails on graphs with negative edge weights
because a later relaxation through a negative edge can decrease an
already-finalised distance.

The running time depends on the priority-queue implementation. With
a binary heap, Dijkstra runs in $O((n + m) \log n)$. With a
Fibonacci heap, the bound improves to $O(m + n \log n)$. For sparse
graphs ($m = O(n)$), both are $O(n \log n)$. For dense graphs, both
are $O(n^{2} \log n)$ in the binary-heap form. In practice, the
binary-heap form is the engineer's default; the Fibonacci-heap
form's constant factors are large enough that the asymptotic gain
rarely repays the implementation cost.

Figure \ref{fig:vol02:ch17:dijkstra-trace} shows the algorithm run
on a five-vertex digraph (the same instance as exercise 17.4). The
relaxation $d(a) = d(b) + 1 = 3$ overwrites the direct-edge
estimate $d(a) = 4$ once $b$ is settled at distance $2$; the
finalisation order is $s, b, a, c, t$. The reference
implementation in
\repopath{volumes/02-form/ch17-discrete-mathematics/code/dijkstra.py}
reproduces the trace and serves as the baseline for the simulation
exercise in section 17.9.

\begin{figure}[ht]
\centering
\begin{tikzpicture}[
  vertex/.style = {circle, draw=engblack, thick, fill=englime!25,
                   inner sep=0pt, minimum size=8mm, font=\small},
  done/.style = {circle, draw=engblue, very thick, fill=engblue!20,
                 inner sep=0pt, minimum size=8mm, font=\small},
  edge/.style = {thick, draw=enggray, ->, >=stealth},
  tedge/.style = {very thick, draw=engblue, ->, >=stealth},
  wlbl/.style = {font=\scriptsize, fill=white, inner sep=1pt},
]

% Five-vertex digraph s -> a, b, c, t used in exercise 17.4.
\node[done]   (s) at (0, 1.5)    {$s$};
\node[vertex] (a) at (2.5, 3)    {$a$};
\node[vertex] (b) at (2.5, 0)    {$b$};
\node[vertex] (c) at (5, 3)      {$c$};
\node[vertex] (t) at (5, 0)      {$t$};

% Edges with weights; tree edges (Dijkstra-found) in engblue.
\draw[tedge] (s) -- node[wlbl] {$4$} (a);
\draw[tedge] (s) -- node[wlbl] {$2$} (b);
\draw[edge]  (a) -- node[wlbl] {$3$} (c);
\draw[edge]  (a) -- node[wlbl] {$1$} (b);
\draw[tedge] (b) -- node[wlbl] {$1$} (a);
\draw[tedge] (b) -- node[wlbl] {$4$} (c);
\draw[tedge] (c) -- node[wlbl] {$2$} (t);
\draw[edge]  (b) -- node[wlbl, pos=0.4] {$5$} (t);

% Final distance labels (the Dijkstra answers).
\node[font=\scriptsize, anchor=south east, engblue]
  at ($(s.north west) + (-0.05, 0.05)$) {$d=0$};
\node[font=\scriptsize, anchor=south, engblue]
  at ($(a.north) + (0, 0.15)$) {$d=3$};
\node[font=\scriptsize, anchor=north, engblue]
  at ($(b.south) + (0, -0.15)$) {$d=2$};
\node[font=\scriptsize, anchor=south, engblue]
  at ($(c.north) + (0, 0.15)$) {$d=6$};
\node[font=\scriptsize, anchor=north, engblue]
  at ($(t.south) + (0, -0.15)$) {$d=8$};

% Finalisation order along the bottom.
\node[font=\footnotesize, anchor=north, enggray]
  at (2.5, -1.4) {finalisation order: $s, b, a, c, t$};
\end{tikzpicture}
\caption{Dijkstra trace on the digraph of exercise 17.4. Tree edges
(thick, engblue) record the predecessor of each settled vertex. The
relaxation $d(a) = d(b) + 1 = 3$ overwrites the initial estimate
$d(a) = 4$ taken from the direct edge $(s, a)$, illustrating how
the algorithm pulls down distances as cheaper routes via earlier
settled vertices appear.}
\label{fig:vol02:ch17:dijkstra-trace}
\end{figure}


\subsection{Shortest path: Bellman-Ford}

Bellman-Ford solves the same single-source shortest-path problem,
under a weaker precondition: the graph may contain negative-weight
edges, provided it contains no negative-weight cycle reachable from
$s$.

\begin{principle}[Bellman-Ford algorithm]
\textbf{Input}: digraph $G = (V, E)$, weights $w$ (possibly
negative), source $s$.
\textbf{Output}: distances $d(s, v)$, or detection of a
negative-weight cycle.

\begin{enumerate}[leftmargin=*]
  \item Initialise as in Dijkstra.
  \item Repeat $n - 1$ times: for every edge $(u, v) \in E$, if
    $d(s, u) + w(u, v) < d(s, v)$, set $d(s, v) = d(s, u) + w(u, v)$.
  \item For every edge $(u, v)$, if $d(s, u) + w(u, v) < d(s, v)$,
    report a negative-weight cycle.
\end{enumerate}
\end{principle}

The running time is $\Theta(nm)$, asymptotically slower than
Dijkstra. The engineer accepts the cost when negative weights are
unavoidable: routing protocols that price links with arbitrage
opportunities, currency-conversion arbitrage detection, and any
problem in which an edge weight encodes a benefit rather than a
cost.

\subsection{Minimum spanning tree: Prim and Kruskal}

A \engterm{spanning tree} of a connected undirected graph is a tree
that includes every vertex. A \engterm{minimum spanning tree} (MST)
of a weighted graph is a spanning tree of minimum total edge
weight. The MST problem is the canonical engineering question of
building the cheapest connected network.

Two algorithms solve MST in efficient working time. Prim's
algorithm grows a single tree from a starting vertex by repeatedly
adding the cheapest edge that connects the tree to a non-tree
vertex. Kruskal's algorithm sorts the edges by weight and adds them
in order, skipping any edge that would create a cycle. Both achieve
$O(m \log n)$ running time. A third algorithm, Bor\r{u}vka's,
predates both and runs in $O(m \log n)$ by growing all components
in parallel and merging on the lightest outgoing edge; it remains
the cleanest MST algorithm to parallelise. Bor\r{u}vka published
the algorithm in 1926 \cite{text:v2ch17-boruvka-1926} to plan an
electrification network for the Moravian region; MST is one of the
few computer-science problems whose first algorithm was driven by
a real engineering deployment.

\begin{principle}[Kruskal's algorithm]
\textbf{Input}: connected weighted graph $G = (V, E)$, weights $w$.
\textbf{Output}: minimum spanning tree $T$.

\begin{enumerate}[leftmargin=*]
  \item Sort the edges in non-decreasing order of weight.
  \item Initialise $T = \emptyset$.
  \item For each edge $(u, v)$ in sorted order: if adding $(u, v)$
    to $T$ does not create a cycle, add it.
  \item Return $T$.
\end{enumerate}
\end{principle}

The cycle test is implemented with a \engterm{union-find} data
structure (also called \engterm{disjoint-set union}) which supports
union and find in nearly constant amortised time, $O(\alpha(n))$
where $\alpha$ is the inverse Ackermann function and is at most $4$
for any input size an engineer will ever face. The dominant cost
of Kruskal is the sort, $O(m \log m) = O(m \log n)$ since $m \leq
n^{2}$.

\begin{principle}[Cut property]
For any partition of the vertex set into a non-empty set $S$ and
its complement $V \setminus S$, the lightest edge crossing the cut
is in some MST. If the lightest crossing edge is unique, it is in
every MST.
\end{principle}

The cut property underwrites both Prim's and Kruskal's algorithms.
Prim's algorithm is the cut property iterated with $S$ as the
current tree's vertex set; Kruskal's is the cut property iterated
across changing cuts as the algorithm builds the tree.

\subsection{Matching: bipartite case}

A \engterm{matching} in a graph $G$ is a set of edges $M \subseteq
E$ no two of which share a vertex. A \engterm{maximum matching} is
a matching of maximum cardinality.

The bipartite case is the engineer's working case. A graph $G = (V,
E)$ is \engterm{bipartite} if its vertex set partitions as $V = A
\cup B$ with every edge going between $A$ and $B$. Bipartite
matchings model job-to-applicant assignments, task-to-machine
schedules, file-to-disk placements, and the assignment problems of
operations research.

The unweighted bipartite matching problem is solved in $O(m
\sqrt{n})$ time by the Hopcroft-Karp algorithm, which alternates
between BFS layering and DFS augmentation. For the weighted case,
the \engterm{Hungarian algorithm} (also called the
Kuhn-Munkres algorithm) finds a maximum-weight perfect matching in
$O(n^{3})$ time. The engineer's working level is recognition: a
problem of the form ``assign $n$ jobs to $n$ machines with cost
$c_{ij}$, minimise total cost'' is the classical assignment problem,
the Hungarian algorithm is the standard solver, and a working
solver is one library call away in any production language.

\begin{archetype}[Optimisation]
A graph algorithm is an optimisation procedure on a discrete
feasible set. The shortest-path problem optimises over walks
between two vertices; the MST problem optimises over spanning
trees; the matching problem optimises over edge subsets satisfying
the matching constraint. Each problem has a standard algorithm with
a known running time, and each algorithm exploits structural
properties of the feasible set (the cut property, the optimal-
substructure property of shortest paths, the augmenting-path
property of matchings) to avoid enumeration. The archetype recurs
in continuous form in Chapter 15 of this volume, in stochastic
form in Volumes IX and X, and in mixed form in operations research
and integer programming.
\end{archetype}

\subsection{Estimation: routing on a continental network}

\begin{estimation}
\textbf{Question.} A continental fibre-optic network has $n
\approx 4 \times 10^{4}$ routers and an average vertex degree
$\bar{d} \approx 5$, sparse and roughly planar in its physical
layout. A network management system computes shortest paths from
each router to every other router once per hour to keep the routing
tables current. Estimate, before reading on, the per-hour
computational cost of the all-pairs shortest-path workload using
Dijkstra from each source with a binary-heap priority queue. State
the working assumption your estimate makes.

\smallskip
\textit{Estimate, before reading on.}

\smallskip
The single-source Dijkstra cost on a graph with $n$ vertices and
$m = \bar{d} n / 2 \approx 10^{5}$ edges is $O((n + m) \log_{2}
n)$. With $n \approx 4 \times 10^{4}$, $\log_{2} n \approx 15$.
The per-source operation count is roughly $(n + m) \log_{2} n
\approx 1.4 \times 10^{5} \cdot 15 \approx 2 \times 10^{6}$.
Multiplied across $n$ sources, the per-hour total is about
$8 \times 10^{10}$ heap operations. A current-generation laptop
performs roughly $10^{10}$ simple operations per second; the
workload takes roughly 8 seconds with a per-op constant of unity,
or roughly 1 to 2 minutes with the per-op constant of a heap
priority-queue (memory allocation, comparison, pointer chasing).
The estimate matters: the workload fits on a single machine with
room to spare. A network operator who proposes a parallel
all-pairs implementation has either misestimated the per-hour cost
or is solving a different problem (such as $k$-shortest paths or
all-pairs with traffic-aware weights).

The same workload on a graph ten times denser ($\bar{d} = 50$, the
regime of a campus-scale Ethernet topology) costs ten times more,
still fits a single machine. On a graph ten times larger ($n = 4
\times 10^{5}$, the continental aggregation of all access-network
nodes), the cost grows by a factor of $100$, pushing into the
single-machine ceiling and motivating a parallel implementation,
landmark-based algorithms (the Galaxy Network architecture; the
A$^{*}$ algorithm with admissible heuristics), or a hierarchical
decomposition.
\end{estimation}

\section{Boolean algebra and logic}\pathtag{standard}

The reader who works with discrete systems will, sooner or later,
have to reason about which inputs cause which outputs in a circuit,
which conditions cause a program to take which branch, or which
combinations of facts cause a constraint to be satisfiable. Boolean
algebra is the apparatus.

\subsection{Boolean operations and truth tables}

A \engterm{Boolean variable} takes values in $\set{0, 1}$, read
\engterm{false} and \engterm{true}. Three operations generate the
algebra: \engterm{conjunction} $\land$ (also AND), \engterm{disjunction}
$\lor$ (also OR), and \engterm{negation} $\lnot$ (also NOT).

A \engterm{truth table} for a Boolean expression in $n$ variables
has $2^{n}$ rows, one for each assignment of values to the
variables, and one column per (sub-)expression. The exponential
growth of the row count in $n$ is the first hint of why exhaustive
analysis fails for large circuits.

\begin{definition}[Boolean operations]
For $a, b \in \set{0, 1}$:
\begin{align*}
a \land b &= 1 \text{ iff } a = b = 1, \\
a \lor b &= 1 \text{ iff } a = 1 \text{ or } b = 1, \\
\lnot a &= 1 - a.
\end{align*}
\end{definition}

Two derived operations are heavily used. The \engterm{exclusive or}
$a \oplus b = (a \lor b) \land \lnot(a \land b)$ is true when
exactly one of $a, b$ is true. The \engterm{implication} $a \implies
b = \lnot a \lor b$ is false only when $a$ is true and $b$ is
false, which is the engineer's working semantics for ``if $a$ then
$b$.''

\subsection{Identities of Boolean algebra}

Boolean algebra obeys the same algebraic discipline as ordinary
algebra, with structural differences worth knowing. Both $\land$
and $\lor$ are commutative and associative. Each distributes over
the other, a stronger property than the integers enjoy: $a \land (b
\lor c) = (a \land b) \lor (a \land c)$ and the dual. The
constants $0$ and $1$ are identities for $\lor$ and $\land$
respectively.

\begin{principle}[De Morgan's laws]
\[
\lnot(a \land b) = (\lnot a) \lor (\lnot b),
\qquad
\lnot(a \lor b) = (\lnot a) \land (\lnot b).
\]
\end{principle}

De Morgan's laws are the working tool by which the engineer rewrites
a Boolean condition into an equivalent form. They are the
standard tool by which a circuit's NAND-only or NOR-only
implementation is derived from an arbitrary Boolean expression,
and they are the standard tool for transforming a complicated
program-branch condition into a form whose negation is easier to
read.

\subsection{Normal forms}

Two normal forms organise Boolean expressions for routine
manipulation.

A \engterm{disjunctive normal form} (DNF) is a disjunction of
conjunctions of literals, where a \engterm{literal} is a variable
or its negation. Every Boolean function on $n$ variables has a DNF:
write the truth table, take one conjunction per row whose output is
$1$, and disjoin them. The DNF derived this way has at most $2^{n}$
conjunctions.

A \engterm{conjunctive normal form} (CNF) is a conjunction of
disjunctions of literals (each disjunction is called a
\engterm{clause}). Every Boolean function has a CNF, derivable by
the dual of the above procedure.

The two forms have different operational uses. DNF is convenient
when one wants to enumerate the truth-true assignments of a
function. CNF is the standard input format of a Boolean
satisfiability solver.

\subsection{SAT}

The \engterm{Boolean satisfiability problem} (SAT) asks: given a
Boolean expression $\varphi$ in $n$ variables, does there exist an
assignment of values to the variables that makes $\varphi$
evaluate to true?

SAT is the canonical NP-complete problem (Cook's theorem, 1971).
For an arbitrary expression in $n$ variables, the worst-case
running time of every known algorithm is exponential in $n$. The
problem is the prototype of computational intractability for the
general case.

The engineer's working level is recognition. SAT solvers (modern
CDCL implementations: MiniSat, Glucose, Kissat, CaDiCaL) routinely
solve industrial instances with millions of variables and tens of
millions of clauses, despite the worst-case bound. Hardware
verification, software model checking, and combinatorial planning
all routinely encode their queries as SAT instances and solve them.
The pattern is general: the worst case is intractable, but the
instances arising from real engineering structure are tractable in
practice. We return to the theme in section 17.8.

Random 3-SAT instances exhibit a striking \engterm{phase
transition} around clause-to-variable ratio $m/n \approx 4.27$
\cite{paper:v2ch17-cheeseman-sat-1991}. For $m/n$ well below this
threshold, almost all instances are easily satisfiable; well above
it, almost all are easily unsatisfiable; right at the threshold,
instances are hard. The reference
\repopath{volumes/02-form/ch17-discrete-mathematics/code/sat_brute.py}
generates instances at the threshold density and measures the
running time of a brute-force solver, the working baseline for the
simulation exercise.

\subsection{Circuit minimisation}

A digital circuit implements a Boolean function as a network of
gates (AND, OR, NOT, NAND, NOR, XOR). The same function can be
implemented by many circuits; the engineer prefers circuits with
fewer gates, fewer levels of logic (lower delay), and lower power
consumption.

Circuit minimisation, the problem of finding a minimum-size circuit
implementing a given function, is computationally hard in the
general case. Karnaugh maps, a graphical procedure for two-level
minimisation in up to four or six variables, were the standard
hand tool of the digital-design generation. Modern logic synthesis
tools (Espresso, ABC, the synthesis pipelines of Synopsys, Cadence,
and Mentor) handle the multi-million-gate scale automatically. The
engineer's working level is to recognise the problem and to read
the synthesis tool's output, not to perform minimisation by hand.

\section{Asymptotic notation}\pathtag{core}

Asymptotic notation expresses how a quantity scales with a
parameter, ignoring constant factors and lower-order terms. The
notation is the engineer's everyday shorthand for the running time
or memory of an algorithm, the error of a numerical method, and the
size of a combinatorial structure as the input grows.

\subsection{The three working bounds}

\begin{definition}[Big-$O$, Big-$\Omega$, Big-$\Theta$]
For functions $f, g: \N \to \R_{>0}$:
\begin{itemize}[leftmargin=*]
  \item $f(n) = O(g(n))$ if there exist constants $c > 0$ and
    $n_{0}$ such that $f(n) \leq c \cdot g(n)$ for all $n \geq
    n_{0}$. We read this ``$f$ is order at most $g$.''
  \item $f(n) = \Omega(g(n))$ if there exist constants $c > 0$ and
    $n_{0}$ such that $f(n) \geq c \cdot g(n)$ for all $n \geq
    n_{0}$. We read this ``$f$ is order at least $g$.''
  \item $f(n) = \Theta(g(n))$ if $f(n) = O(g(n))$ and $f(n) =
    \Omega(g(n))$. We read this ``$f$ is of order exactly $g$.''
\end{itemize}
\end{definition}

The notation is asymmetric in the equality symbol: $f(n) = O(g(n))$
is read ``$f \in O(g)$'' rather than ``$f$ equals $O(g)$.'' The
symbol is, in this context, set membership in the equivalence class
of functions bounded above by a constant multiple of $g$. The
engineer's macro $\oom{n^{2}}$ in this book renders this set
membership.

\subsection{The working hierarchy}

A short list of growth classes covers nearly every engineering
running time. From slowest-growing to fastest:
\[
O(1) \subset O(\log \log n) \subset O(\log n) \subset O(\log^{k} n)
\subset O(n^{\epsilon}) \subset O(n) \subset O(n \log n)
\]
\[
\subset O(n^{2}) \subset O(n^{3}) \subset O(n^{k}) \subset O(c^{n})
\subset O(n!) \subset O(n^{n}).
\]
The hierarchy is the engineer's first scaling literacy. An
algorithm that is $O(n^{2})$ on inputs of size $n = 10^{4}$
performs $10^{8}$ operations, comfortable on a current laptop. The
same algorithm on $n = 10^{6}$ performs $10^{12}$ operations,
overnight on a server. On $n = 10^{8}$, $10^{16}$ operations,
beyond a single-machine deadline. Quadratic running time is the
boundary between feasible and infeasible at modern scale; an
engineer who proposes a quadratic algorithm for a million-element
input has misread the scale.

Figure \ref{fig:vol02:ch17:growth-classes} plots the classes
on a semilog axis up to $n = 64$. The single-machine
operations budget for a one-year run (current as of 2026, roughly
$3 \times 10^{17}$ operations) sits as a horizontal line across the
plot. The class crossings against this line set the engineer's
feasibility frontier: linear and $n \log n$ remain comfortably
below, quadratic crosses near $n \approx 10^{9}$, cubic near
$n \approx 10^{6}$, exponential $2^{n}$ near $n \approx 60$, and
factorial near $n \approx 21$. The reference table at
\repopath{volumes/02-form/ch17-discrete-mathematics/data/feasibility_frontier.csv}
records the per-class frontier for single-machine and modest-cluster
budgets.

\begin{figure}[p]
\centering
\begin{tikzpicture}
\begin{semilogyaxis}[
  width = 14cm,
  height = 16cm,
  xlabel = {input size $n$},
  ylabel = {operations},
  xmin = 1, xmax = 64,
  ymin = 1, ymax = 1e22,
  grid = both,
  major grid style = {gray!30},
  minor grid style = {gray!15},
  legend pos = north west,
  legend cell align = {left},
  legend style = {font=\small},
  title = {Asymptotic growth classes on a semilog plot},
  every axis title/.style = {below right, at={(0.02, 0.98)}, font=\small\bfseries},
]

% log n
\addplot[domain=1:64, samples=80, very thick, draw=engblue]
  {ln(x)/ln(2)};
\addlegendentry{$\log_{2} n$};

% n
\addplot[domain=1:64, samples=80, very thick, draw=engblack, dashed]
  {x};
\addlegendentry{$n$};

% n log n
\addplot[domain=1:64, samples=80, very thick, draw=enggray]
  {x * ln(x)/ln(2)};
\addlegendentry{$n \log_{2} n$};

% n^2
\addplot[domain=1:64, samples=80, very thick, draw=engamber]
  {x^2};
\addlegendentry{$n^{2}$};

% n^3
\addplot[domain=1:64, samples=80, very thick, draw=engamber, dashed]
  {x^3};
\addlegendentry{$n^{3}$};

% 2^n
\addplot[domain=1:64, samples=80, very thick, draw=engred]
  {2^x};
\addlegendentry{$2^{n}$};

% n!
\addplot[domain=1:21, samples=22, very thick, draw=engred, dashed]
  {factorial(x)};
\addlegendentry{$n!$};

% Single-machine budget (10^18) horizontal line.
\addplot[domain=1:64, samples=2, very thick, draw=engblue, dotted]
  {1e18};
\addlegendentry{single-machine budget ($\approx 10^{18}$)};

\end{semilogyaxis}
\end{tikzpicture}
\caption{The asymptotic hierarchy on a semilog plot, $1 \leq n \leq
64$. Logarithmic, linear, $n \log n$, and quadratic classes remain
below the single-machine operations budget ($\approx 10^{18}$
operations per year on a current-generation laptop, current as of
2026) across the full range. Cubic crosses the budget near $n
\approx 10^{6}$. Exponential $2^{n}$ crosses near $n \approx 60$.
Factorial $n!$ crosses near $n \approx 21$. The figure is the
chapter's central pedagogical picture: the engineer reads off the
ceiling for an enumeration before launching the search. The
hierarchy is the reading discipline of sections 17.5 and 17.8.}
\label{fig:vol02:ch17:growth-classes}
\end{figure}


\subsection{Working examples}

A linear scan over $n$ items is $\Theta(n)$. Sorting $n$ items by
the standard comparison-based methods (mergesort, heapsort,
introsort) is $\Theta(n \log n)$. Binary search over a sorted array
of $n$ items is $\Theta(\log n)$. The classical $n \times n$ matrix
multiplication is $\Theta(n^{3})$, with Strassen's algorithm
$\Theta(n^{\log_{2} 7}) = \Theta(n^{2.807\ldots})$ as a working
mastery-level alternative. Solving a dense $n \times n$ linear
system by Gaussian elimination is $\Theta(n^{3})$. Brute-force
enumeration over all subsets of an $n$-element set is $\Theta(2^{n})$.
Brute-force enumeration over all permutations is $\Theta(n!)$.

\subsection{The constant matters in practice}

Asymptotic analysis ignores the constant factor. The engineer does
not. An $O(n)$ algorithm with a constant factor of $10^{4}$ is
beaten by an $O(n^{2})$ algorithm with a constant factor of $1$ for
$n < 10^{4}$. In the working regime, constant factors and lower-order
terms can dominate the asymptotic class. The engineer's discipline
is to do both an asymptotic analysis (will the algorithm scale?)
and a concrete benchmark on representative inputs (does it run in
acceptable time on the inputs we care about?).

\begin{mastery}
A second discipline of asymptotic analysis is the choice of cost
model. The standard model is the \engterm{RAM model}: a constant-time
read or write of a constant-size word. The model is appropriate when
arithmetic involves machine-precision numbers and the algorithm's
data fits comfortably in memory. When numbers are arbitrary-precision
integers (cryptography, number theory), or when the data exceeds
memory and disk I/O dominates (external-memory algorithms), the cost
model changes and the asymptotic class can shift. CLRS Chapter 1
\cite{text:cormen-clrs} discusses the model assumptions explicitly;
the engineer who uses asymptotic analysis on a non-RAM problem
without checking the model is at risk of a one-line objection from a
reviewer.
\end{mastery}

\section{Recurrences and generating functions}\pathtag{standard}

Recurrences arise whenever a problem of size $n$ is reduced to one
or more subproblems of smaller size. Divide-and-conquer algorithms
(mergesort, quicksort, FFT, Strassen) admit a recurrence for their
running time; recursive data structures (binary trees, balanced
BSTs, heaps) admit recurrences for their height; and many counting
problems admit a recurrence for their answer.

\subsection{Solving a recurrence}

A \engterm{recurrence relation} expresses a sequence in terms of
its earlier values. The Fibonacci recurrence $F_{n} = F_{n-1} +
F_{n-2}$ with $F_{0} = 0, F_{1} = 1$ is the canonical example.

Three working techniques solve recurrences in routine engineering
work.

\textbf{Substitution} guesses the form of the solution and verifies
it by induction. For the recurrence $T(n) = 2 T(n/2) + n$ with
$T(1) = 1$, the guess $T(n) = O(n \log n)$ verifies: assume $T(k)
\leq c k \log_{2} k$ for all $k < n$, substitute, and obtain
$T(n) \leq 2 (c (n/2) \log_{2}(n/2)) + n = c n \log_{2} n - c n + n
\leq c n \log_{2} n$ for $c \geq 1$.

\textbf{Recursion-tree} expansion writes the recurrence's recursion
tree and sums the work at each level. For $T(n) = 2 T(n/2) + n$ the
tree has $\log_{2} n$ levels, each performing $\Theta(n)$ work,
giving $\Theta(n \log n)$. Figure \ref{fig:vol02:ch17:recursion-tree}
shows the tree explicitly. The technique is the engineer's fastest
path to the right answer when the recurrence has a clean structure.

\begin{figure}[ht]
\centering
\begin{tikzpicture}[
  every node/.style = {font=\footnotesize},
  node distance = 0.8cm and 0.4cm,
  level distance = 1.05cm,
  level 1/.style={sibling distance=8cm},
  level 2/.style={sibling distance=4cm},
  level 3/.style={sibling distance=2cm},
  level 4/.style={sibling distance=1cm},
  cost/.style = {rectangle, draw=engblack, fill=englime!20, thick,
                 inner sep=2pt, minimum width=1.2cm},
]
% Recursion tree for T(n) = 2 T(n/2) + n.
\node[cost] (root) {$n$}
  child {node[cost] {$n/2$}
    child {node[cost] {$n/4$}
      child {node[cost] {$n/8$} edge from parent[draw=enggray]}
      child {node[cost] {$n/8$} edge from parent[draw=enggray]}
      edge from parent[draw=enggray]}
    child {node[cost] {$n/4$}
      child {node[cost] {$n/8$} edge from parent[draw=enggray]}
      child {node[cost] {$n/8$} edge from parent[draw=enggray]}
      edge from parent[draw=enggray]}
    edge from parent[draw=enggray]}
  child {node[cost] {$n/2$}
    child {node[cost] {$n/4$}
      child {node[cost] {$n/8$} edge from parent[draw=enggray]}
      child {node[cost] {$n/8$} edge from parent[draw=enggray]}
      edge from parent[draw=enggray]}
    child {node[cost] {$n/4$}
      child {node[cost] {$n/8$} edge from parent[draw=enggray]}
      child {node[cost] {$n/8$} edge from parent[draw=enggray]}
      edge from parent[draw=enggray]}
    edge from parent[draw=enggray]};

% Level-cost annotations on the right.
\node[anchor=west, engblue] at (7.5, 0)     {level cost $= n$};
\node[anchor=west, engblue] at (7.5, -1.05) {$2 \cdot (n/2) = n$};
\node[anchor=west, engblue] at (7.5, -2.10) {$4 \cdot (n/4) = n$};
\node[anchor=west, engblue] at (7.5, -3.15) {$8 \cdot (n/8) = n$};

% Depth axis on the left.
\node[anchor=east] at (-7.5, 0)     {depth $0$};
\node[anchor=east] at (-7.5, -1.05) {depth $1$};
\node[anchor=east] at (-7.5, -2.10) {depth $2$};
\node[anchor=east] at (-7.5, -3.15) {depth $3$};

\node[anchor=west, font=\small] at (-7.5, -4.0)
  {Total depth $\log_{2} n$, total work $\Theta(n \log n)$.};
\end{tikzpicture}
\caption{Recursion tree for $T(n) = 2 T(n/2) + n$, the running-time
recurrence for mergesort. At each level the work sums to $n$; the
tree has $\log_{2} n$ levels (one base-$2$ halving per level until
the subproblems hit the base case); the total work is the level
count times the per-level work, $\Theta(n \log n)$. The same
picture, with the level work replaced by $\Theta(n^{\log_{b} a})$,
underwrites case 2 of the master theorem.}
\label{fig:vol02:ch17:recursion-tree}
\end{figure}


\textbf{Master theorem} packages the recursion-tree analysis for a
common class of divide-and-conquer recurrences.

\begin{principle}[Master theorem]
For a recurrence
\[
T(n) = a T(n/b) + f(n),
\]
with $a \geq 1, b > 1$, $T$ defined on positive integers, and $f$
asymptotically positive, the asymptotic behaviour is determined by
$f(n)$ relative to $n^{\log_{b} a}$:
\begin{itemize}[leftmargin=*]
  \item If $f(n) = O(n^{\log_{b} a - \epsilon})$ for some $\epsilon
    > 0$, then $T(n) = \Theta(n^{\log_{b} a})$.
  \item If $f(n) = \Theta(n^{\log_{b} a})$, then $T(n) =
    \Theta(n^{\log_{b} a} \log n)$.
  \item If $f(n) = \Omega(n^{\log_{b} a + \epsilon})$ for some
    $\epsilon > 0$ and a regularity condition holds, then $T(n) =
    \Theta(f(n))$.
\end{itemize}
\end{principle}

The three cases distinguish the regime in which the recursion's
work is dominated by the leaves (case 1), evenly distributed
across levels (case 2), or dominated by the root (case 3). The
mergesort recurrence $T(n) = 2 T(n/2) + n$ is case 2, giving
$\Theta(n \log n)$. The binary-search recurrence $T(n) = T(n/2) +
1$ is case 2 with $\log_{b} a = 0$, giving $\Theta(\log n)$.
Karatsuba's $T(n) = 3 T(n/2) + n$ is case 1, giving $\Theta(n^{\log_{2}
3}) = \Theta(n^{1.585\ldots})$. Strassen's $T(n) = 7 T(n/2) +
n^{2}$ is case 1, giving $\Theta(n^{\log_{2} 7})$.

\subsection{Generating functions}

A \engterm{generating function} encodes a sequence $(a_{0}, a_{1},
a_{2}, \ldots)$ as a formal power series
\[
A(x) = \sum_{n \geq 0} a_{n} x^{n}.
\]
The variable $x$ is formal: convergence is rarely the engineer's
concern. Generating functions convert sequence operations
(shifting, convolution, weighted summation) into algebraic
operations on power series, often turning a recurrence into an
algebraic equation that solves in closed form.

The Fibonacci generating function $F(x) = \sum_{n \geq 0} F_{n}
x^{n}$ satisfies $F(x) = x F(x) + x^{2} F(x) + x$, which solves to
$F(x) = x / (1 - x - x^{2})$. Partial-fraction decomposition then
delivers the closed form
\[
F_{n} = \frac{1}{\sqrt{5}} \left( \phi^{n} - \psi^{n} \right),
\quad \phi = \frac{1 + \sqrt{5}}{2}, \psi = \frac{1 - \sqrt{5}}{2}.
\]

The engineer's working level for generating functions is
recognition. When a recurrence resists substitution and master
theorem, generating functions are the standard next tool, and the
machinery exists in symbolic-computation packages (Mathematica's
\texttt{RSolve}, SageMath, SymPy) that handle the routine cases
without manual calculation.

\section{Number theory engineers actually use}\pathtag{standard}

Number theory in this section means the small fragment that
recurs in cryptography, coding theory, hashing, and the working
internals of computer arithmetic.

\subsection{Modular arithmetic}

Two integers $a, b$ are \engterm{congruent modulo $n$}, written $a
\equiv b \pmod{n}$, if $n$ divides $a - b$. Modular arithmetic
defines $+, -, \cdot$ on the residue classes $\set{0, 1, \ldots, n
- 1}$ in the natural way: $(a + b) \bmod n$, $(a \cdot b) \bmod n$,
and so on.

The set $\Z/n\Z = \set{0, 1, \ldots, n-1}$ with these operations is
a \engterm{commutative ring}. It is a \engterm{field} (every
non-zero element has a multiplicative inverse) if and only if $n$
is prime. The structural difference is the source of the
cancellation rule's failure for non-prime moduli, the example we
gave in Chapter 1.

\subsection{Greatest common divisor and the Euclidean algorithm}

\begin{definition}[Greatest common divisor]
The \engterm{greatest common divisor} of two integers $a, b$, not
both zero, written $\gcd(a, b)$, is the largest positive integer
that divides both. Two integers with $\gcd(a, b) = 1$ are
\engterm{coprime}.
\end{definition}

The \engterm{Euclidean algorithm} computes $\gcd(a, b)$ by repeated
remainder.

\begin{principle}[Euclidean algorithm]
\textbf{Input}: non-negative integers $a, b$, not both zero.
\textbf{Output}: $\gcd(a, b)$.

\begin{enumerate}[leftmargin=*]
  \item If $b = 0$, return $a$.
  \item Otherwise, return $\gcd(b, a \bmod b)$.
\end{enumerate}
\end{principle}

The algorithm terminates because each step strictly decreases the
second argument. The number of steps is $O(\log(\min(a, b)))$, by
a Fibonacci-style argument: the worst case is when consecutive
arguments are consecutive Fibonacci numbers, in which case the
ratio of successive arguments is the golden ratio $\phi$ and the
algorithm takes $\log_{\phi} a$ steps.

The \engterm{extended Euclidean algorithm} returns, in addition to
$\gcd(a, b)$, integers $x, y$ such that $ax + by = \gcd(a, b)$.
B\'{e}zout's identity guarantees the existence of such $x, y$. The
extended algorithm is the standard tool for computing modular
inverses: when $\gcd(a, n) = 1$, $ax + ny = 1$ implies $a x \equiv 1
\pmod{n}$, that is, $x \equiv a^{-1} \pmod{n}$.

\subsection{Primes}

A \engterm{prime} is an integer $p \geq 2$ whose only positive
divisors are $1$ and $p$. The fundamental theorem of arithmetic
states that every positive integer factors uniquely into primes (up
to the order of factors).

The prime-counting function $\pi(n)$ counts primes at most $n$.
The prime number theorem states $\pi(n) \sim n / \ln n$, that is,
$\pi(n) / (n / \ln n) \to 1$ as $n \to \infty$. For working
estimates, the density of primes near $n$ is roughly $1 / \ln n$;
near $n = 10^{18}$, about one in every $41$ integers is prime.

Two algorithmic facts of working importance:

\textbf{Primality testing} of an integer $n$ is in deterministic
polynomial time (the AKS algorithm,
\cite{paper:v2ch17-aks-2004}), but the working choice is a
probabilistic Miller-Rabin test, which runs in $O(k \log^{3} n)$
for $k$ rounds and has failure probability at most $4^{-k}$ on
composite inputs. For $k = 40$ rounds, the failure probability is
below $10^{-24}$, well beyond cryptographic confidence. With a
fixed witness set of the first thirteen primes, the test is
deterministic for $n < 3.3 \times 10^{24}$ (Sorenson and Webster,
2017), which covers nearly every working application short of
RSA-key generation. The reference
\repopath{volumes/02-form/ch17-discrete-mathematics/code/miller_rabin.py}
implements the test in this form and finds the first $100$ primes
above $10^{18}$ in a fraction of a second on a current laptop.

\textbf{Integer factorisation} is, as far as we know, hard: the
best classical algorithm (the general number-field sieve) runs in
sub-exponential time $\exp(O((\log n)^{1/3} (\log \log n)^{2/3}))$,
which is impractical for $n$ above a few hundred bits. Shor's
quantum algorithm runs in polynomial time on a sufficiently large
fault-tolerant quantum computer, which is, current as of 2026, an
engineering programme of unsettled timeline. The hardness of
factorisation is the basis of RSA cryptography (\S below).

\subsection{Cryptographic primitives, recognition level}

Three primitives recur often enough that an engineer should
recognise them.

\textbf{RSA}. Choose two large primes $p, q$, compute $n = pq$ and
the Euler totient $\varphi(n) = (p - 1)(q - 1)$. Choose an exponent
$e$ coprime to $\varphi(n)$ and compute its modular inverse $d
\equiv e^{-1} \pmod{\varphi(n)}$. The public key is $(n, e)$; the
private key is $d$. Encryption is $c = m^{e} \bmod n$; decryption
is $m = c^{d} \bmod n$. Correctness rests on Fermat's little
theorem and its Euler generalisation.

\textbf{Diffie-Hellman}. Two parties agree on a prime $p$ and a
generator $g$ of the multiplicative group $(\Z/p\Z)^{\ast}$. Each
party picks a private exponent and exchanges $g^{a} \bmod p$ and
$g^{b} \bmod p$. The shared secret is $g^{ab} \bmod p$. The
security rests on the hardness of the discrete logarithm problem
in the multiplicative group.

\textbf{Hash functions}. A cryptographic hash function maps
arbitrary-length strings to fixed-length digests, with the working
properties of preimage resistance, second-preimage resistance, and
collision resistance. SHA-256, specified in NIST FIPS 180-4
\cite{std:v2ch17-nist-fips-180-4}, and BLAKE3 are the engineer's
working choices, current as of 2026. SHA-1, deprecated by NIST in
2011 and broken in practice by the SHAttered collision attack of
2017, illustrates the half-life of a cryptographic primitive: the
hash that was the workhorse of TLS through the 2000s is no longer
acceptable for collision-resistance applications, and an engineer
who deploys SHA-1 today for new code has ignored the trajectory.

The engineer's working level is recognition rather than design. A
working engineer never designs a cryptographic primitive from
scratch: the field has produced enough subtle disasters that the
discipline is to use a vetted library implementation and to keep
the cryptographic boundary thin. The reader is referred to Volume
VII for the depth treatment.

\section{Failure: combinatorial explosion in design spaces}\pathtag{core}

The chapter's failure mode is the discovery that a search over a
combinatorial design space is infeasible.

A \engterm{combinatorial explosion} is the rapid growth of the
size of a discrete structure with the size of its parameters. The
hierarchy is the one we met in section 17.5: polynomial growth is
dominated by exponential, exponential by factorial, factorial by
double exponential. The failure mode is the engineer who proposes
to enumerate the structure in a regime where enumeration is
impossible.

\subsection{The arithmetic of impossibility}

A current-generation laptop performs roughly $10^{10}$ simple
operations per second, current as of 2026. A workstation cluster of
$10^{3}$ nodes performs roughly $10^{13}$ operations per second. A
year of continuous run is $\approx 3 \times 10^{7}$ seconds. The
working ceiling for a single-machine search is therefore $\approx 3
\times 10^{17}$ operations; for a cluster, $3 \times 10^{20}$.

The corresponding feasibility frontier is around $n = 60$ for
$2^{n}$ (single machine) and $n = 70$ for $2^{n}$ (cluster). For
$n!$, the frontier is around $n = 20$ (single machine) and $n = 22$
(cluster). For $(n!)^{n}$ or $2^{n^{2}}$, the frontier is in the
single digits.

A design space whose size exceeds the feasibility frontier cannot
be searched exhaustively. The engineer's response is one of three:

\begin{enumerate}[leftmargin=*]
  \item Reformulate the problem so the search space is smaller.
    Often the actual problem has structure (locality, decomposition,
    monotonicity) the original formulation hides. Exposing the
    structure shrinks the search by orders of magnitude. The
    transformation of a generic constraint problem into a SAT
    instance, then solving with a CDCL solver that exploits
    learning, is the routine modern instance.
  \item Use an exact algorithm that is exponential in the worst
    case but practical on the structured instances that arise. The
    SAT, MIP, and CP solver communities have built such algorithms.
    Their worst-case bound is unchanged; their average-case and
    structured-case performance is dramatically better.
  \item Settle for an approximate or heuristic solution. Greedy
    algorithms, local search, simulated annealing, evolutionary
    algorithms, and modern reinforcement-learning policies all
    sacrifice the guarantee of optimality for tractable running
    time. The engineer's discipline is to characterise the
    approximation: how far from optimal, with what probability, on
    which input distribution.
\end{enumerate}

\subsection{The same archetype, met in continuous form}

The reader who has seen Chapter 16 (numerical methods) and will see
Chapter 18 (mathematical modelling) will recognise the failure as a
combinatorial form of the \engterm{curse of dimensionality}: the
phenomenon by which the size of a discretised continuous space
grows exponentially with its dimension. A continuous optimisation
in $d$ dimensions, discretised at $k$ points per axis, requires
$k^{d}$ evaluations; for $k = 10$ and $d = 20$, that is $10^{20}$
evaluations, beyond the cluster ceiling.

The combinatorial and the continuous forms are the same archetype.
The engineer who has internalised the scaling literacy of section
17.5 reads the curse before encountering its formal statement.

\subsection{The named failures}

Three classes of engineering failure trace to a misjudged
combinatorial estimate.

\textbf{Brute-force key search.} An engineer who underestimates the
key-space size of a cryptographic primitive can deploy a system
whose security is breakable. The DES cipher, with its $56$-bit key,
was a working cipher in 1977 and a broken cipher by the late 1990s,
when distributed brute-force key search became feasible at $2^{56}
\approx 7 \times 10^{16}$ operations. The lesson is not that DES
was poorly designed for its time; it is that an engineer who
deploys a system in 2026 with a key space below $2^{80}$ has
ignored the trajectory.

\textbf{Test coverage.} A software system with $n$ Boolean
configuration flags has $2^{n}$ configurations. For $n = 30$, the
configuration space is $\approx 10^{9}$, beyond exhaustive
testing. The engineer who reports ``we tested every configuration''
on a $30$-flag system is either using combinatorial test design
(pairwise covering arrays) or is reporting in error. Pairwise
covering arrays, for example, cover every pair of flag values in
$O(n \log n)$ tests: a small constant-factor expansion that
catches a documented majority of multi-flag interaction bugs.

\textbf{Generative design searches.} A modern generative-design
tool that searches over $10^{30}$ candidate geometries cannot
enumerate. Either the tool exploits a learned surrogate (a neural
network that approximates the cost function and shrinks the
effective search), or it accepts a heuristic, or it reformulates
the problem. An engineer who relies on the tool's output without
understanding its search strategy has no honest account of why
the answer is good.

\begin{principle}[Combinatorial-explosion check]
Before committing to a discrete search, estimate the search-space
size. If the estimate exceeds the operations-budget ceiling
($\approx 10^{18}$ for a single machine, $\approx 10^{20}$ for a
modest cluster, current as of 2026), the search must either
exploit structure, accept approximation, or be reformulated. The
engineer who does not perform the estimate is committed to one of
those three responses by default; performing it lets them choose
which.
\end{principle}

\section{Project}\pathtag{core}

\begin{project}[A real network as a graph]
\textbf{Question.} Take a real network of your choice and produce a
graph model that other engineers could query.

\textbf{Network choice}. Pick one of the following or one of
comparable scale:
\begin{itemize}[leftmargin=*]
  \item A city transit network (subway or bus): nodes are stations,
    edges are direct service segments.
  \item An electrical-grid section (a transmission or distribution
    grid for a region): nodes are buses or substations, edges are
    lines.
  \item A water-distribution network: nodes are junctions or
    reservoirs, edges are pipes.
  \item A social network drawn from a public dataset: nodes are
    accounts, edges are follows or friendships.
  \item A supply chain or logistics network: nodes are facilities,
    edges are shipping routes.
  \item A software-dependency graph: nodes are packages, edges are
    dependencies.
  \item A biological network (protein-protein interaction, neural
    connectome): nodes are entities, edges are interactions.
\end{itemize}

\textbf{Documentation.} Write up the data source (URL, DOI, or
citation), the abstraction choices (what counts as a node, what
counts as an edge, what is left out, why), and the resulting graph's
basic statistics ($n$, $m$, average degree, connected-component
count).

\textbf{Computations.} Compute, for the largest connected component:
\begin{enumerate}[leftmargin=*]
  \item The diameter, by all-pairs shortest path or BFS from each
    vertex.
  \item The global clustering coefficient.
  \item The degree distribution. Plot it on log-log axes. Estimate
    a power-law exponent if the distribution looks heavy-tailed; a
    Poisson approximation if it does not.
  \item The five highest-degree vertices and the five
    highest-betweenness vertices. Identify what each of these
    corresponds to in the underlying physical or social system.
\end{enumerate}

\textbf{Discussion.} In a 1500-word reflection, address:
\begin{itemize}[leftmargin=*]
  \item What the graph metrics reveal about the network's
    resilience: which vertices, if removed, would fragment
    connectivity, and which vertices' failure would be invisible.
  \item How the abstraction choices affected the conclusions:
    treating bus and subway as a single multimodal graph versus
    two coupled graphs would change the numbers; including or
    excluding suburban branches would change them.
  \item One conclusion the engineer would draw if asked for a
    one-sentence summary, and the principal caveat the engineer
    would attach.
\end{itemize}

\textbf{Tools.} Paper, pencil, optionally a programming
environment with graph-algorithm primitives (NetworkX, igraph,
graph-tool, or equivalent). The standard library calls are
documented; the work is in the abstraction choices and the
discussion.

\textbf{Deliverable.} A 10-page document with the data source, the
graph statistics, the plots, the identified hubs, and the
reflection. Repository path:
\repopath{volumes/02-form/ch17-discrete-mathematics/projects/}.

\textbf{Time.} Two to three hours of data acquisition and cleaning,
two to three hours of graph computation, two hours of writing.
The total work fits in a long weekend.

\textbf{Phases.}
\begin{enumerate}[leftmargin=*]
  \item \textit{Acquisition} (2--3 h). Find a documented source for
    the network. For transit, the General Transit Feed Specification
    (GTFS) feeds publish stop lists and stop-time tables. For power
    grids, regional grid operators (ENTSO-E in Europe, the U.S.
    Energy Information Administration) publish substation and
    line-segment data with deliberate aggregation. For software
    dependencies, the registries (PyPI, npm, Cargo) expose
    dependency metadata. Pick a source whose abstractions you can
    name and whose licence permits the work.
  \item \textit{Abstraction} (1 h). Write the node and edge
    definitions explicitly. Record the omissions (one-way bus
    segments, transformers, optional dependencies) and the
    rationale.
  \item \textit{Computation} (2--3 h). Implement the four
    computations (diameter, clustering, degree distribution, hubs)
    using NetworkX, igraph, or graph-tool. Verify on a small
    hand-checked subgraph before running on the full data.
  \item \textit{Interpretation} (1 h). Identify the high-degree and
    high-betweenness vertices and locate them in the underlying
    physical or social system. The interpretation is the part the
    library cannot do.
  \item \textit{Write-up} (2 h). The 1500-word reflection plus the
    plots.
\end{enumerate}

\textbf{Success criteria.}
\begin{enumerate}[leftmargin=*]
  \item The data source resolves to a citation (URL, DOI) and the
    abstraction choices are stated explicitly enough that a
    different engineer could reproduce the graph.
  \item The four computations are reported with their algorithm
    name, parameter choices, and observed running time.
  \item The degree distribution plot shows axes, units, and a
    fitted curve (power-law or Poisson) with the fit's parameters.
  \item The high-degree vertices are identified by their real-world
    name (Times Square, Substation X, the \texttt{lodash} package)
    and the engineer offers a one-sentence explanation of why each
    is a hub.
  \item The reflection takes a position on the network's resilience
    and identifies the principal caveat the engineer would attach.
\end{enumerate}

\textbf{Common failure modes.}
\begin{enumerate}[leftmargin=*]
  \item \textit{Treating a disconnected graph as connected.}
    Diameter is undefined on disconnected graphs; compute it on the
    largest connected component, not on the raw input.
  \item \textit{Loading a multigraph as a simple graph.} Many real
    datasets carry parallel edges (two bus routes on the same
    segment); silently collapsing them changes the degree
    distribution. Decide explicitly and document.
  \item \textit{Computing all-pairs shortest paths on a million-node
    graph.} The cost scales as $\Theta(n \cdot (n + m) \log n)$
    with Dijkstra from each source. A road-network with $n = 10^{6}$
    saturates a single machine; use a sample of source vertices for
    the diameter estimate or a landmark-based algorithm.
  \item \textit{Fitting a power law without checking the bound.} A
    log-log plot looks linear over many distributions that are not
    power laws. The reader is referred to Clauset, Shalizi, and
    Newman's hypothesis-testing procedure for the proper fit.
  \item \textit{Confusing degree with betweenness.} The two
    centrality measures agree on canonical hubs but disagree on
    bottlenecks: a low-degree vertex on the only path between two
    large communities has high betweenness and low degree. Report
    both.
\end{enumerate}
\end{project}

\section{Exercises}

The exercises mix calculation, derivation, estimation, simulation,
design, diagnosis, failure analysis, and judgement per the
editorial settlement on exercise types. Full solutions are provided
for the calculation and derivation exercises; the estimation
exercises receive published reference estimates with the editor's
working; the design, diagnosis, failure-analysis, and judgement
exercises receive rubrics rather than single answers.

\subsection*{Calculation}\pathtag{core}

\begin{exercise}
Compute the number of distinct ways to arrange the letters of
ENGINEERING. State the multinomial coefficient that governs the
count.
\end{exercise}

\paragraph{Solution sketch.} The word has 11 letters with class
sizes E:3, N:3, I:2, G:2, R:1. The multinomial coefficient is
\[
\binom{11}{3, 3, 2, 2, 1} = \frac{11!}{3!\, 3!\, 2!\, 2!\, 1!}
= \frac{39\,916\,800}{6 \cdot 6 \cdot 2 \cdot 2 \cdot 1}
= \frac{39\,916\,800}{144} = 277\,200.
\]

\begin{exercise}
Compute $\binom{20}{7}$ exactly, then verify that
$\binom{20}{7} = \binom{19}{6} + \binom{19}{7}$ using Pascal's
identity.
\end{exercise}

\paragraph{Solution sketch.}
$\binom{20}{7} = 20! / (7!\, 13!) = 77\,520$. Pascal:
$\binom{19}{6} = 27\,132$ and $\binom{19}{7} = 50\,388$; sum is
$27\,132 + 50\,388 = 77\,520$, agreeing with $\binom{20}{7}$.

\begin{exercise}
For the graph $G$ on $V = \set{1, 2, 3, 4, 5, 6}$ with edges $E =
\set{\set{1,2}, \set{2,3}, \set{3,4}, \set{4,5}, \set{5,6},
\set{1,6}, \set{2,5}}$, compute (a) the degree sequence, (b) the
diameter, (c) the global clustering coefficient.
\end{exercise}

\paragraph{Solution sketch.}
(a) Degrees: $\deg(1) = 2$ (neighbours $2, 6$), $\deg(2) = 3$
(neighbours $1, 3, 5$), $\deg(3) = 2$ (neighbours $2, 4$),
$\deg(4) = 2$ (neighbours $3, 5$), $\deg(5) = 3$ (neighbours $2,
4, 6$), $\deg(6) = 2$ (neighbours $1, 5$). Degree sequence
$(3, 3, 2, 2, 2, 2)$. Handshake check: $\sum \deg = 14 = 2 m$
with $m = 7$.

(b) The graph is the 6-cycle $1\text{-}2\text{-}3\text{-}4\text{-}5\text{-}6\text{-}1$
with one chord $\set{2, 5}$. Distances:
$d(1, 4) = 3$ via $1\text{-}2\text{-}3\text{-}4$ or
$1\text{-}6\text{-}5\text{-}4$, and $d(1, 4) = 3$ via the chord
($1\text{-}2\text{-}5\text{-}4$) does not improve. All other
pairs at distance $\leq 3$; diameter is $3$.

(c) Triangles: $\set{2, 5}$ is a chord; we check whether any
triple $\set{i, j, k}$ with all three pairwise edges exists. The
only candidate triangles involve the chord: $\set{2, 5, 1}$
($\set{2,5}, \set{1,2}$ yes, $\set{1,5}$ no), $\set{2, 5, 6}$
($\set{5,6}$ yes, $\set{2,6}$ no), $\set{2, 5, 3}$ ($\set{2,3}$
yes, $\set{3,5}$ no), $\set{2, 5, 4}$ ($\set{4,5}$ yes, $\set{2,4}$
no). No triangles; clustering coefficient is $0$ at every vertex
with $\deg \geq 2$, and the global clustering coefficient is $0$.

\begin{exercise}
On the directed weighted graph with vertices $\set{s, a, b, c, t}$
and edges $(s, a, 4)$, $(s, b, 2)$, $(a, c, 3)$, $(a, b, 1)$, $(b,
c, 4)$, $(b, a, 1)$, $(c, t, 2)$, $(b, t, 5)$, run Dijkstra from
$s$ and report $d(s, v)$ for every $v$. Indicate the order in which
vertices are finalised.
\end{exercise}

\paragraph{Solution sketch.}
The trace is shown in figure
\ref{fig:vol02:ch17:dijkstra-trace}. Initial: $d(s) = 0$, all
others $\infty$. Finalise $s$, relax: $d(a) = 4$, $d(b) = 2$.
Finalise $b$ (smallest unvisited at $2$); relax: $d(a) =
\min(4, 2 + 1) = 3$, $d(c) = 2 + 4 = 6$, $d(t) = 2 + 5 = 7$.
Finalise $a$ ($d = 3$); relax: $d(c) = \min(6, 3 + 3) = 6$
(no change). Finalise $c$ ($d = 6$); relax: $d(t) =
\min(7, 6 + 2) = 8$. Finalise $t$. Final distances:
$d(s, s) = 0$, $d(s, b) = 2$, $d(s, a) = 3$, $d(s, c) = 6$,
$d(s, t) = 8$. Finalisation order: $s, b, a, c, t$.

\begin{exercise}
On the undirected weighted graph with vertices $\set{1, 2, 3, 4,
5}$ and edges $\set{1,2}: 1$, $\set{2,3}: 4$, $\set{1,3}: 3$,
$\set{3,4}: 2$, $\set{4,5}: 5$, $\set{2,4}: 6$, run Kruskal's
algorithm and report the MST and its total weight.
\end{exercise}

\paragraph{Solution sketch.}
Sort edges by weight: $\set{1,2}{:}1$, $\set{3,4}{:}2$,
$\set{1,3}{:}3$, $\set{2,3}{:}4$, $\set{4,5}{:}5$, $\set{2,4}{:}6$.
Process in order. Add $\set{1,2}$ (no cycle); add $\set{3,4}$ (no
cycle); add $\set{1,3}$ (no cycle, joins the $\set{1,2}$ and
$\set{3,4}$ components); skip $\set{2,3}$ (would create cycle
$1\text{-}2\text{-}3$); add $\set{4,5}$ (extends tree to vertex 5);
$\set{2,4}$ would create a cycle but the tree already has $n - 1 = 4$
edges, so the algorithm halts. MST: $\set{1,2}, \set{3,4},
\set{1,3}, \set{4,5}$, total weight $1 + 2 + 3 + 5 = 11$. The
reference implementation in
\repopath{volumes/02-form/ch17-discrete-mathematics/code/kruskal.py}
reproduces this.

\begin{exercise}
Compute the truth table of $f(a, b, c) = (a \land b) \lor (\lnot
a \land c)$ in all eight rows. Write a CNF and a DNF for $f$.
\end{exercise}

\paragraph{Solution sketch.}
Truth table (rows in $abc$ order, $0$ to $7$): $f(0,0,0) = 0$;
$f(0,0,1) = 1$; $f(0,1,0) = 0$; $f(0,1,1) = 1$; $f(1,0,0) = 0$;
$f(1,0,1) = 0$; $f(1,1,0) = 1$; $f(1,1,1) = 1$. Truth-true rows:
$(0,0,1), (0,1,1), (1,1,0), (1,1,1)$.

DNF (one conjunction per truth-true row):
\[
f = (\lnot a \land \lnot b \land c) \lor (\lnot a \land b \land c)
  \lor (a \land b \land \lnot c) \lor (a \land b \land c).
\]
Simplification by combining minterms gives the original form
$f = (a \land b) \lor (\lnot a \land c)$.

CNF (one clause per truth-false row, with each literal negated):
$f(0,0,0) = 0$ gives clause $(a \lor b \lor c)$;
$f(0,1,0) = 0$ gives $(a \lor \lnot b \lor c)$;
$f(1,0,0) = 0$ gives $(\lnot a \lor b \lor c)$;
$f(1,0,1) = 0$ gives $(\lnot a \lor b \lor \lnot c)$. CNF is the
conjunction of these four clauses. Simplification gives
$(a \lor c) \land (\lnot a \lor b)$.

\begin{exercise}
Solve the recurrence $T(n) = 4 T(n/2) + n^{2}$ with $T(1) = 1$ by
master theorem. State which case applies and the resulting
$\Theta(\cdot)$ class.
\end{exercise}

\paragraph{Solution sketch.}
Identify $a = 4$, $b = 2$, $f(n) = n^{2}$. Compute
$n^{\log_{b} a} = n^{\log_{2} 4} = n^{2}$. Compare $f(n) =
\Theta(n^{2}) = \Theta(n^{\log_{b} a})$: this is case 2 of the
master theorem. The conclusion is $T(n) = \Theta(n^{\log_{b} a}
\log n) = \Theta(n^{2} \log n)$. The recurrence is the running
time of the standard square-matrix multiplication algorithm
expressed in divide-and-conquer form.

\begin{exercise}
Compute $\gcd(2024, 1547)$ by the Euclidean algorithm, showing each
remainder step. Then run the extended Euclidean algorithm and
report integers $x, y$ such that $2024 x + 1547 y = \gcd(2024,
1547)$.
\end{exercise}

\paragraph{Solution sketch.}
Euclidean steps:
$2024 = 1 \cdot 1547 + 477$;
$1547 = 3 \cdot 477 + 116$;
$477 = 4 \cdot 116 + 13$;
$116 = 8 \cdot 13 + 12$;
$13 = 1 \cdot 12 + 1$;
$12 = 12 \cdot 1 + 0$. The gcd is $1$, so $2024$ and $1547$ are
coprime.

Back-substitute for the extended algorithm. From $1 = 13 - 1 \cdot
12$ and $12 = 116 - 8 \cdot 13$: $1 = 13 - (116 - 8 \cdot 13) = 9
\cdot 13 - 116$. From $13 = 477 - 4 \cdot 116$: $1 = 9(477 - 4
\cdot 116) - 116 = 9 \cdot 477 - 37 \cdot 116$. From $116 = 1547 -
3 \cdot 477$: $1 = 9 \cdot 477 - 37(1547 - 3 \cdot 477) = 120
\cdot 477 - 37 \cdot 1547$. From $477 = 2024 - 1547$: $1 = 120
(2024 - 1547) - 37 \cdot 1547 = 120 \cdot 2024 - 157 \cdot 1547$.
Check: $120 \cdot 2024 - 157 \cdot 1547 = 242\,880 - 242\,879 = 1$.
So $x = 120, y = -157$ satisfy $2024 x + 1547 y = 1$.

\subsection*{Derivation}\pathtag{standard}

\begin{exercise}
Prove the handshake lemma $\sum_{v \in V} \deg(v) = 2 \abs{E}$ for
any undirected graph by a direct double count of edge-endpoint
incidences.
\end{exercise}

\paragraph{Solution sketch.}
Define an \engterm{incidence} to be a pair $(v, e)$ with $v \in V$,
$e \in E$, and $v \in e$. Count incidences two ways. Fix a vertex
$v$: the number of incidences with first coordinate $v$ is exactly
$\deg(v)$, since $\deg(v)$ counts the edges containing $v$. Summing
over $v$ gives $\sum_{v} \deg(v)$ incidences in total. Fix an edge
$e = \set{u, v}$: the number of incidences with second coordinate
$e$ is exactly $2$, namely $(u, e)$ and $(v, e)$. Summing over $e$
gives $2 \abs{E}$ incidences in total. The two counts are of the
same set, so they are equal. The corollary follows: $\sum
\deg(v)$ is even, so the number of odd-degree vertices is even.

\begin{exercise}
Prove that a tree on $n \geq 1$ vertices has exactly $n - 1$
edges, by induction on $n$.
\end{exercise}

\paragraph{Solution sketch.}
Induct on $n$. Base case $n = 1$: a one-vertex tree has $0$ edges,
and $n - 1 = 0$. Inductive step: assume every tree on $k$ vertices
has $k - 1$ edges, $1 \leq k < n$. Let $T$ be a tree on $n$
vertices. A tree has at least one leaf (a vertex of degree $1$);
otherwise every vertex has degree $\geq 2$, the sum of degrees is
$\geq 2n > 2(n - 1)$, contradicting the handshake lemma plus the
cycle-free property $\abs{E} \leq n - 1$. Pick a leaf $v$ and
remove it together with its incident edge. The result $T - v$ has
$n - 1$ vertices, is still connected (removing a leaf cannot
disconnect a tree), and is still cycle-free; so $T - v$ is a tree
on $n - 1$ vertices with $(n - 1) - 1 = n - 2$ edges by induction.
Adding back $v$ and its incident edge gives $n - 1$ edges total.

\begin{exercise}
Prove the cut property of MST: for any partition $V = S \cup (V
\setminus S)$ with both sides non-empty, the lightest edge crossing
the cut is in some MST. Use an exchange argument: take any MST $T$
and show that if a lightest cut edge $e$ is not in $T$, swapping
$e$ for some edge $e'$ in $T$ that crosses the cut produces a tree
of equal or smaller weight.
\end{exercise}

\paragraph{Solution sketch.}
Let $e = \set{u, v}$ with $u \in S$, $v \in V \setminus S$ be a
lightest edge crossing the cut. Let $T$ be any MST. If $e \in T$,
we are done. Otherwise add $e$ to $T$. The result $T + e$ has $n$
edges and contains a unique cycle $C$ (the cycle is unique because
$T$ was a spanning tree). The cycle $C$ enters $S$ via $e$ and
must leave $S$ along some other edge $e'$ (since the endpoints of
$C$ at $u$ and $v$ alternate sides of the cut at least twice as we
trace $C$); $e'$ crosses the cut and $e'$ is in $T$. Now form $T'
= T + e - e'$: $T'$ is connected (we have removed an edge that
sits on a cycle in $T + e$), has $n - 1$ edges, and so is a
spanning tree. Its weight is $w(T) + w(e) - w(e')$. Since $e$
is a lightest cut edge, $w(e) \leq w(e')$, so $w(T') \leq w(T)$.
Since $T$ was an MST, $w(T') = w(T)$ and $T'$ is also an MST, and
$T'$ contains $e$. So some MST contains $e$.

\begin{exercise}
Prove de Morgan's law $\lnot(a \land b) = \lnot a \lor \lnot b$ by
a four-row truth table.
\end{exercise}

\paragraph{Solution sketch.}
The four rows of $(a, b)$ give:
\begin{center}
\begin{tabular}{cccc}
$a$ & $b$ & $\lnot(a \land b)$ & $\lnot a \lor \lnot b$ \\
\hline
0 & 0 & 1 & 1 \\
0 & 1 & 1 & 1 \\
1 & 0 & 1 & 1 \\
1 & 1 & 0 & 0 \\
\end{tabular}
\end{center}
Columns agree row by row, so the two expressions are
identically equal.

\begin{exercise}
Prove that every Boolean function on $n$ variables admits a DNF
representation. Construct the DNF directly from the truth table.
State the upper bound on the number of conjunctions.
\end{exercise}

\paragraph{Solution sketch.}
Let $f: \set{0, 1}^{n} \to \set{0, 1}$. For each input $\vec{x} =
(x_{1}, \ldots, x_{n}) \in \set{0, 1}^{n}$ with $f(\vec{x}) = 1$,
form the conjunction
\[
m_{\vec{x}} = \ell_{1} \land \ell_{2} \land \cdots \land \ell_{n},
\quad
\ell_{i} = \begin{cases} x_{i} & \text{if } x_{i} = 1, \\
                          \lnot x_{i} & \text{if } x_{i} = 0.
            \end{cases}
\]
Each $m_{\vec{x}}$ is true on exactly the one assignment $\vec{x}$
and false on all others. The disjunction $\bigvee_{\vec{x} :
f(\vec{x}) = 1} m_{\vec{x}}$ is true on exactly the assignments
where $f$ is true and so equals $f$. The number of conjunctions is
the number of truth-true rows, at most $2^{n}$.

\begin{exercise}
Prove the master theorem case 2: if $f(n) = \Theta(n^{\log_{b} a})$
in the recurrence $T(n) = a T(n/b) + f(n)$, then $T(n) =
\Theta(n^{\log_{b} a} \log n)$. Use the recursion-tree argument and
sum the work at each level.
\end{exercise}

\paragraph{Solution sketch.}
The recursion tree for $T(n) = a T(n/b) + f(n)$ has the root
performing $f(n)$ work and $a$ children each on subproblems of
size $n/b$. At depth $k$ the tree has $a^{k}$ subproblems each of
size $n / b^{k}$. The work at depth $k$ is $a^{k} \cdot f(n /
b^{k})$. Using $f(n) = c_{1} n^{\log_{b} a}$ at the top and
$c_{2} n^{\log_{b} a}$ at the bottom (the $\Theta$ bound),
substitute:
\[
a^{k} f(n / b^{k}) = c \cdot a^{k} (n / b^{k})^{\log_{b} a}
= c \cdot a^{k} n^{\log_{b} a} / (b^{\log_{b} a})^{k}
= c \cdot a^{k} n^{\log_{b} a} / a^{k}
= c \cdot n^{\log_{b} a}.
\]
The work at each level is $\Theta(n^{\log_{b} a})$, independent
of $k$. The tree has $\log_{b} n$ levels (the recursion halves
subproblem size by factor $b$ each step, until size drops to a
constant). Total work is $\Theta(n^{\log_{b} a}) \cdot \log_{b} n
= \Theta(n^{\log_{b} a} \log n)$.

\begin{exercise}
Prove that the Euclidean algorithm terminates in $O(\log \min(a,
b))$ steps. (Hint: show that two consecutive remainders satisfy
$r_{i+2} < r_{i} / 2$.)
\end{exercise}

\paragraph{Solution sketch.}
The Euclidean algorithm produces a remainder sequence $r_{0} = a$,
$r_{1} = b$, and $r_{i+1} = r_{i-1} \bmod r_{i}$ for $i \geq 1$,
strictly decreasing and non-negative. We show $r_{i+2} <
r_{i} / 2$. By definition, $r_{i+2} = r_{i} \bmod r_{i+1}$. Case 1:
if $r_{i+1} \leq r_{i} / 2$, then $r_{i+2} < r_{i+1} \leq r_{i} /
2$. Case 2: if $r_{i+1} > r_{i} / 2$, then $r_{i} = 1 \cdot
r_{i+1} + r_{i+2}$, so $r_{i+2} = r_{i} - r_{i+1} < r_{i} - r_{i} /
2 = r_{i} / 2$. Either way, $r_{i+2} < r_{i} / 2$. So every two
steps halve the remainder; after $2 \log_{2} \min(a, b)$ steps the
remainder drops below $1$, terminating the algorithm. Step count
is $O(\log \min(a, b))$.

\begin{exercise}
Prove Fermat's little theorem: for a prime $p$ and an integer $a$
not divisible by $p$, $a^{p-1} \equiv 1 \pmod{p}$. (Hint: the map
$x \mapsto a x \bmod p$ permutes $\set{1, \ldots, p-1}$; multiply
both sides of the resulting equation by inverses.)
\end{exercise}

\paragraph{Solution sketch.}
The set $S = \set{1, 2, \ldots, p - 1}$ is the multiplicative
group of $\Z/p\Z$. The map $\mu_{a}: x \mapsto a x \bmod p$ is a
bijection on $S$: if $a x \equiv a y \pmod{p}$ then $a (x - y)
\equiv 0$, and since $\gcd(a, p) = 1$, $p \mid x - y$, so $x = y$
in $S$. So $\mu_{a}$ permutes $S$. The product over $S$ equals
the product over $\mu_{a}(S)$:
\[
\prod_{x \in S} x = \prod_{x \in S} (a x) = a^{p - 1}
\prod_{x \in S} x \pmod{p}.
\]
Cancel $\prod_{x \in S} x = (p - 1)!$ (which is coprime to $p$)
from both sides to get $1 \equiv a^{p - 1} \pmod{p}$.

\subsection*{Estimation}\pathtag{core}

\begin{exercise}
Estimate the number of distinct shuffles of a $52$-card deck.
Express the result in the closest power of $10$. Estimate the
ratio of this count to the number of atoms in the observable
universe ($\approx 10^{80}$).
\end{exercise}

\paragraph{Solution sketch.}
$52! = 8.07 \times 10^{67}$. Closest power of $10$ is $10^{68}$.
Ratio to atoms: $10^{68} / 10^{80} = 10^{-12}$, that is, the deck
count is about one trillionth of the atom count. The reference
estimate: every shuffle of a well-randomised deck has, with
overwhelming probability, never been seen before.

\begin{exercise}
A graph algorithm runs in $O(n^{2})$ on an input of size $n =
10^{4}$ in $5$ seconds. Estimate, in seconds and in human-readable
form (minutes, hours, days, years), the time the same algorithm
would take on inputs of size $n = 10^{6}$, $n = 10^{8}$, and $n =
10^{10}$. State the assumption your estimate makes.
\end{exercise}

\paragraph{Solution sketch.}
Assume the asymptotic class is the running time at the relevant
scale (the constant factor is unchanged and lower-order terms are
negligible). At $n = 10^{4}$, runtime is $5$ s; the algorithm
performs roughly $(10^{4})^{2} = 10^{8}$ basic operations per
$5$ s, that is, $2 \times 10^{7}$ ops/s. Multiplying inputs by
$10^{2}$ multiplies time by $10^{4}$ ($O(n^{2})$):
$n = 10^{6}$: $5 \times 10^{4}$ s $\approx 14$ hours;
$n = 10^{8}$: $5 \times 10^{8}$ s $\approx 16$ years;
$n = 10^{10}$: $5 \times 10^{12}$ s $\approx 1.6 \times 10^{5}$
years. The principal assumption that fails in practice is the
memory bound: $n = 10^{8}$ already requires terabyte-scale storage
for an adjacency matrix, and the running time is dominated by
memory bandwidth long before $10^{10}$.

\begin{exercise}
Estimate the number of operations a current laptop, a $10^{3}$-node
cluster, and a $10^{6}$-node supercomputer can each perform in one
year. Use these numbers to identify, for each platform, the largest
$n$ for which a $\Theta(2^{n})$ enumeration is feasible.
\end{exercise}

\paragraph{Solution sketch.}
One year $= 3.15 \times 10^{7}$ s. Per-platform ops/s, current as
of 2026: laptop $\approx 10^{10}$, cluster ($10^{3}$ nodes)
$\approx 10^{13}$, supercomputer ($10^{6}$ nodes) $\approx 10^{16}$.
Per-year operations: laptop $3 \times 10^{17}$, cluster $3 \times
10^{20}$, supercomputer $3 \times 10^{23}$. Solve $2^{n} \leq B$
for each ceiling $B$: laptop $n \leq \log_{2}(3 \times 10^{17})
\approx 58$; cluster $n \leq \log_{2}(3 \times 10^{20}) \approx
68$; supercomputer $n \leq \log_{2}(3 \times 10^{23}) \approx 78$.
The exercise's headline: a thousandfold speedup buys ten extra
bits.

\begin{exercise}
A road network in a small city has $\approx 5000$ intersections
and $\approx 12000$ road segments. Estimate the total storage cost
of an adjacency-matrix representation versus an adjacency-list
representation. State which representation you would choose and
why.
\end{exercise}

\paragraph{Solution sketch.}
Adjacency matrix: $n \times n$ bits with $n = 5000$ gives $2.5
\times 10^{7}$ entries; at one byte each (a boolean stored in a
byte, the common case in a dense-matrix library), about $25$ MB,
or at one float each (for weighted edges), about $100$ MB.
Adjacency list: $2 m + n$ entries (each edge listed at each
endpoint plus a per-vertex header), with $m = 12\,000$ that is
$\approx 30\,000$ integer references, about $0.24$ MB at four
bytes per integer. The adjacency list wins by two orders of
magnitude. Choose adjacency list: the graph is sparse (average
degree $\bar{d} = 2 m / n = 4.8$, far from the dense regime where
$\bar{d} \approx n$), and shortest-path and BFS workloads scan
edges per vertex, the operation the list supports in $\Theta(\deg
v)$.

\begin{exercise}
A deck of $52$ playing cards is shuffled. Estimate the probability
that the result places the four aces consecutively (in any order,
in any of the $49$ possible starting positions). Compare with a
direct calculation.
\end{exercise}

\paragraph{Solution sketch.}
Estimate: pretend the aces are placed uniformly at random in the
52 positions. The first ace lands somewhere; the second has to
land adjacent to it, probability roughly $2/51$; the third has to
extend the block, roughly $2/50$; the fourth, roughly $2/49$.
Product is $\approx (2/51)(2/50)(2/49) \approx 6 \times 10^{-5}$,
of order $10^{-4}$.

Direct calculation: count favourable arrangements as
$49$ starting positions $\times 4!$ orderings of the aces $\times
48!$ orderings of the other cards, total $49 \cdot 4! \cdot 48!$.
Total arrangements $52!$. Probability is $49 \cdot 4! \cdot 48! /
52! = 49 \cdot 24 / (52 \cdot 51 \cdot 50 \cdot 49) = 24 / (52
\cdot 51 \cdot 50) = 24 / 132\,600 \approx 1.81 \times 10^{-4}$.
Estimate and calculation agree to a factor of three.

\subsection*{Simulation}\pathtag{mastery}

\begin{exercise}
Implement Dijkstra's algorithm in a programming environment of
your choice, with a binary-heap priority queue. Run it on a
randomly-generated graph with $n = 10^{4}$ vertices and $m = 5 n$
edges and verify the running time empirically against the
theoretical $O((n + m) \log n)$ bound by repeating with $n \in
\set{10^{3}, 10^{4}, 10^{5}}$ and plotting time versus
$(n + m) \log n$.
\end{exercise}

\paragraph{Solution sketch.}
The reference implementation lives at
\repopath{volumes/02-form/ch17-discrete-mathematics/code/dijkstra.py}.
Generate $G(n, m)$ with $m = 5n$ by sampling without replacement
from the $n(n-1)/2$ unordered pairs and assigning uniform weights
on $[1, 100]$. For each $n$, run Dijkstra from $10$ randomly
chosen sources, average. Tabulate $T(n)$ versus $(n + m) \log_{2}
n$; the points should lie close to a line through the origin in a
linear-scale plot. Expect deviations at the low end (overhead
dominates) and at the high end (cache effects appear). A linear
fit's slope is the per-heap-op constant, typically tens to
hundreds of nanoseconds on a current laptop.

\begin{exercise}
Implement Kruskal's algorithm with a union-find data structure on
a randomly-generated graph and confirm it produces an MST of the
same total weight as Prim's algorithm. Use $n = 10^{3}$ and $m = 4
n$ as the working scale.
\end{exercise}

\paragraph{Solution sketch.}
Reference Kruskal at
\repopath{volumes/02-form/ch17-discrete-mathematics/code/kruskal.py}.
Implement Prim's algorithm with a min-heap keyed on the cheapest
known edge into the tree. Run both on the same random graph (seed
the RNG); the edge sets differ when weights tie, but the total
weights agree exactly. On $n = 10^{3}$, $m = 4 \times 10^{3}$,
Kruskal runs in about $1$ ms on a current laptop. Repeating on
$n \in \set{10^{3}, 10^{4}, 10^{5}}$ confirms $O(m \log n)$
scaling.

\begin{exercise}
Implement a brute-force SAT solver that tries all $2^{n}$
assignments. Run it on randomly-generated $3$-CNF instances with
$n \in \set{10, 15, 20, 25, 30}$ and report the running time.
Plot time on a log axis against $n$.
\end{exercise}

\paragraph{Solution sketch.}
Reference at
\repopath{volumes/02-form/ch17-discrete-mathematics/code/sat_brute.py}.
Use clause-to-variable ratio $m/n \approx 4.25$ to land near the
phase-transition density. Expected wall-clock on a current laptop:
$n = 10$ in microseconds; $n = 20$ in a few seconds; $n = 25$ in
roughly $2$ minutes; $n = 30$ in roughly $1$ hour (the time
doubles per added variable, as $2^{n}$ predicts). Plotting
$\log_{10} T$ versus $n$ gives a straight line with slope
$\log_{10} 2 \approx 0.30$. A CDCL solver on the same instances
finishes in milliseconds across the full range; the comparison is
the chapter's empirical illustration of the worst-case versus
working-case gap.

\begin{exercise}
Generate an Erd\H{o}s-R\'{e}nyi random graph $G(n, p)$ for $n =
1000$ and $p \in \set{1/n, 2/n, 4/n, 8/n}$. Compute the size of
the largest connected component for each $p$ and plot it. Identify
the threshold $p \approx 1/n$ at which a giant component emerges.
\end{exercise}

\paragraph{Solution sketch.}
Reference at
\repopath{volumes/02-form/ch17-discrete-mathematics/code/erdos_renyi.py}.
For $n = 1000$ and average degree $c = np$, the giant-component
fraction $\rho$ solves $\rho = 1 - e^{-c \rho}$. Expected values:
$c = 0.5$: $\rho \approx 0$ (no giant component);
$c = 1$: $\rho \approx 0$ (right at the threshold; the giant is
$o(n)$ in size);
$c = 2$: $\rho \approx 0.80$;
$c = 4$: $\rho \approx 0.98$;
$c = 8$: $\rho \approx 1$. The empirical largest-component sizes
should match these fractions to within tens of vertices. The plot
of $|\text{LCC}| / n$ versus $c$ traces out the classical Erdos-Renyi
phase transition.

\begin{exercise}
Implement the Miller-Rabin primality test and use it to find the
first $100$ primes greater than $10^{18}$. Report the running time.
Compare with a deterministic trial-division test on a
representative input and report the speedup.
\end{exercise}

\paragraph{Solution sketch.}
Reference at
\repopath{volumes/02-form/ch17-discrete-mathematics/code/miller_rabin.py}.
Expected running time on a current laptop: the search for $100$
primes above $10^{18}$ takes roughly $0.1$ s (the prime density
near $10^{18}$ is $\approx 1/\ln(10^{18}) \approx 1/41$, so the
search examines roughly $4000$ candidates; each Miller-Rabin test
runs in microseconds at this size). Trial division on a candidate
of order $10^{18}$ would scan up to $\sqrt{10^{18}} = 10^{9}$
divisors, taking roughly $1$ second per candidate; the
$4000$-candidate search would take more than an hour. Speedup is
about four orders of magnitude.

\subsection*{Design}\pathtag{standard}

\begin{exercise}
Design an adjacency-list representation for a directed weighted
graph that supports the following operations efficiently: enumerate
all out-neighbours of a vertex, enumerate all in-neighbours of a
vertex, look up the weight of a given edge, add a new edge, remove
an existing edge. State the asymptotic cost of each operation in
your design and identify the operation whose cost is the worst.
\end{exercise}

\paragraph{Discussion.}
A working design carries two per-vertex hash maps: $\mathrm{out}[u]$
maps a head vertex $v$ to the edge weight $w(u, v)$, and
$\mathrm{in}[v]$ symmetrically maps a tail vertex $u$ to the same
weight. Enumeration of out-neighbours of $u$ is $\Theta(\deg^{+}
u)$ by iterating $\mathrm{out}[u]$; enumeration of in-neighbours
of $v$ is $\Theta(\deg^{-} v)$. Edge-weight lookup
$(u, v) \mapsto w$ is $\Theta(1)$ amortised via $\mathrm{out}[u]$.
Edge insertion is $\Theta(1)$ amortised (two hash-table inserts).
Edge deletion is $\Theta(1)$ amortised (two hash-table deletes).
The duplication doubles storage to $\Theta(n + m)$ at the cost of
$\Theta(n + m)$ extra. The worst operation in this design is
out-neighbour enumeration when the implementation is locality-naive
(hash maps fragment cache lines); for a static graph used by
shortest-path benchmarks, a compressed sparse row (CSR)
representation with sorted neighbour arrays gives the same
asymptotics with order-of-magnitude better constants.

\begin{exercise}
Design a graph schema for a public transit network that you would
hand to a working software engineer. The schema must include
station identifiers, line identifiers, segment weights (time and
distance), service-frequency information, and accessibility
metadata. Defend two of your design choices against an imagined
alternative.
\end{exercise}

\paragraph{Discussion.}
A workable schema has node tables (station id, name, geographic
coordinates, accessibility flags), line tables (line id, mode,
operator), and edge tables (station id pair, line id, distance,
travel time, headway, accessibility flags). Two defensible choices:
\textit{(i)} represent each line-on-segment as a separate parallel
edge rather than collapsing parallel services into a single
multi-modal edge; the choice supports per-line frequency, per-line
incident reporting, and routing that respects transfers, at the
cost of inflating the edge count. Defending against the
alternative: the routing algorithm needs to know which line is
being taken (transfer penalties depend on it), so the line
information cannot be discarded; representing it on edges rather
than via a side-channel keeps the routing query a single pass over
the graph. \textit{(ii)} carry segment travel time as the primary
edge weight rather than distance; users care about time, not
distance, and the routing query reduces to ordinary shortest path.
Defending against the alternative: a distance-weighted routing
would correctly find geographically short routes that miss
infrequent services and detours through slow modes.

\begin{exercise}
Design a test plan for a Boolean expression in $10$ variables that
arises in a safety-critical setting. The test plan must justify
its coverage strategy against the $2^{10} = 1024$ assignments. State
how a pairwise covering array would compare with a full enumeration.
\end{exercise}

\paragraph{Discussion.}
$2^{10} = 1024$ is small enough to enumerate, and the safety-critical
context tips the balance toward exhaustive testing: a single
covered-pair miss is a costly residual risk. A defensible test plan
states the coverage criterion (every assignment of the $10$
variables exercised at least once), the test harness's run time
(at $1$ ms per evaluation, the suite takes about $1$ second), and
the failure-reporting strategy (which assignments produced
unexpected outputs, traced back to the equation's structure). A
pairwise covering array would cover every pair of variable values
in approximately $10$--$20$ tests (the standard table for $2^{10}$
gives $13$ tests), a $50$-fold reduction; the cost is that
three-way interactions are not guaranteed covered. For
safety-critical code, the savings do not justify the lost coverage;
for non-safety-critical code at the same scale, the pairwise array
is the working choice. The crossover is dictated by the
consequence of a missed interaction, not by the absolute test count.

\subsection*{Diagnosis}\pathtag{standard}

\begin{exercise}
A junior engineer reports that Dijkstra's algorithm gave a
shortest-path distance of $-3$ on a graph that contained an edge of
weight $-1$. Diagnose what went wrong, name the precondition the
engineer violated, and state which algorithm should have been used.
\end{exercise}

\paragraph{Solution sketch.}
The precondition Dijkstra requires is non-negative edge weights.
The presence of any edge of weight $-1$ violates the invariant that
the first time the algorithm finalises a vertex $u$, $d(s, u)$ is
the true shortest-path distance: a later relaxation across the
negative edge can decrease the distance below the value the
algorithm fixed. The reported $-3$ is either an artefact of the
silent violation (some implementations return a partial answer
that the user reads as final) or a downstream computation that
trusted the wrong distance. The correct algorithm for graphs with
negative edges and no negative-weight cycle is Bellman-Ford,
$O(nm)$; if negative cycles may be present, Bellman-Ford detects
them in its final pass.

\begin{exercise}
A graph-algorithm benchmark measures running time of $10$ seconds on
$n = 10^{3}$ and $40$ seconds on $n = 2 \times 10^{3}$. State which
asymptotic class the benchmark is consistent with and propose two
follow-up measurements that would distinguish $\Theta(n^{2})$ from
$\Theta(n \log n)$ definitively.
\end{exercise}

\paragraph{Solution sketch.}
Doubling $n$ multiplied the time by $4$. $\Theta(n^{2})$ predicts a
factor of $4$ exactly, consistent. $\Theta(n \log n)$ predicts a
factor of $2 \cdot (\log_{2} 2000 / \log_{2} 1000) = 2 \cdot
(10.97 / 9.97) \approx 2.20$, well below the observed $4$. The
benchmark is consistent with $\Theta(n^{2})$ and not with
$\Theta(n \log n)$. Two follow-ups that strengthen the
discrimination: \textit{(i)} measure at $n = 10^{4}$: $n^{2}$
predicts $1000$ s, $n \log n$ predicts $26.6$ s; the gap is two
orders of magnitude. \textit{(ii)} plot $T(n) / (n \log n)$ on a
log scale across $n \in \set{10^{3}, 2 \times 10^{3}, 5 \times
10^{3}, 10^{4}}$; a constant ratio confirms $n \log n$, a linearly
rising ratio confirms $n^{2}$.

\begin{exercise}
A junior engineer presents a Boolean simplification claiming
$(a \lor b) \land \lnot a = b$. Diagnose the error precisely, state
the corrected identity, and identify the de Morgan or
distributivity step the engineer omitted.
\end{exercise}

\paragraph{Solution sketch.}
Distribute $\lnot a$ over the disjunction: $(a \lor b) \land
\lnot a = (a \land \lnot a) \lor (b \land \lnot a) = 0 \lor (b
\land \lnot a) = b \land \lnot a$. The corrected identity is
$(a \lor b) \land \lnot a = b \land \lnot a$, not $b$. The error:
the simplification dropped the $\land \lnot a$ on $b$. The
counterexample is $a = 1, b = 1$: left side $(1 \lor 1) \land 0 =
0$, the claimed right side $b = 1$. The engineer's correct chain
uses distributivity ($\land$ over $\lor$) plus the complement law
$a \land \lnot a = 0$; the missed step is the failure to retain
the $\lnot a$ factor on the surviving $b$ term.

\begin{exercise}
A modular-arithmetic routine reports that the inverse of $6$ modulo
$15$ is $11$. Verify or refute. If the inverse does not exist, state
the precondition the routine should have checked and identify the
correct response.
\end{exercise}

\paragraph{Solution sketch.}
The claim is refuted: $6 \cdot 11 = 66 = 4 \cdot 15 + 6 \equiv 6
\pmod{15}$, not $1$. More structurally, $\gcd(6, 15) = 3 \neq 1$,
so $6$ has no multiplicative inverse modulo $15$: the precondition
for $a^{-1} \bmod n$ to exist is $\gcd(a, n) = 1$. The routine
should have called the extended Euclidean algorithm, which returns
the gcd; if the gcd exceeds $1$, the routine should refuse the
operation (raise an exception, return a sentinel value, or
otherwise signal the precondition failure) rather than return any
number that satisfies neither the inverse nor a sensible default.

\subsection*{Failure analysis}\pathtag{core}

\begin{exercise}
The DES cipher used a $56$-bit key and was the U.S. data-encryption
standard from 1977 to 2001. Argue, in 200 to 300 words, that the
cipher's eventual deprecation was a predictable consequence of the
combinatorial-explosion estimate of section 17.8 plus Moore's-law
trajectory, and identify which of the three engineer's responses
to combinatorial explosion was selected by the cryptographic
community when they replaced DES with AES.
\end{exercise}

\paragraph{Solution sketch.}
In 1977 the DES key space $2^{56} \approx 7.2 \times 10^{16}$
exceeded a single-machine search budget by roughly four orders of
magnitude (the typical 1977 mainframe ran at $\sim 10^{6}$
operations per second, giving $\sim 10^{13}$ ops/year). Moore's
trajectory predicted that single-machine throughput would roughly
double every two years, shrinking the security margin by one bit
every two years. Predicted year of feasibility: $1977 + 2 \times 13
\approx 2003$, which matched the EFF DES Cracker (1998) and
distributed DESCHALL effort (1997) within a few years. The
deprecation was a predictable arithmetic consequence of the
combinatorial-explosion estimate against the Moore trajectory.

The community's response was the first of the three: reformulate
the search by enlarging the key. AES (Rijndael, 2001) supports
$128$-, $192$-, and $256$-bit keys; even at sustained Moore's-law
growth, $2^{128} \approx 3.4 \times 10^{38}$ stays beyond the
feasibility frontier of any plausible classical computer for a
century. The community did not accept an approximate or
probabilistic answer (cryptography requires worst-case guarantees)
nor did it exploit structure in the cipher (the structure-exploiting
attack on DES, differential cryptanalysis, was already known when
AES was designed; AES was built to resist it). Reformulation by
enlargement was the only durable answer.

\begin{exercise}
A software product with $n = 25$ Boolean configuration flags ships
with the claim that ``every configuration has been tested.'' Argue,
in 200 to 300 words, that the claim is either using combinatorial
test design or is misreporting. Estimate the cost of pairwise
covering-array testing as an alternative and state under what
condition you would accept the claim at face value.
\end{exercise}

\paragraph{Solution sketch.}
$2^{25} = 33\,554\,432$ configurations. At $1$ second per
configuration test, exhaustive testing takes more than a year on a
single machine. At $1$ ms per test, $9$ hours, plausible but
expensive; at $1$ microsecond per test (an unrealistic budget for
anything but a pure-CPU smoke check), $33$ seconds. The claim
``every configuration has been tested'' is either using a per-test
cost below microsecond scale (suggesting a trivially shallow test,
in which case the claim is technically true but informationally
vacuous) or it is misreporting. The likely truth is that the
vendor ran tests on a small selected subset and reported the
optimistic phrase.

A pairwise covering array for $n = 25$ Boolean flags requires
between $10$ and $20$ test cases (the precise count from
constructions of Cohen et al. is $13$ or $14$). At $1$ second per
test, the suite runs in $20$ seconds and covers every pair of
flag values, catching the documented majority of two-flag
interaction bugs. Higher-order covering arrays scale slowly: a
three-way cover for $25$ flags is on the order of $30$--$50$ tests.

The claim is acceptable at face value only when \textit{(i)} the
per-test cost is genuinely under a microsecond and the test
exercises a measured-significant code path, or \textit{(ii)} the
domain admits structural reduction (most flags are independent and
the conjoint behaviour is provably the same as the per-flag
behaviour). Both conditions are unusual; the working presumption
is that the claim is misreported.

\begin{exercise}
A road-network routing application using Dijkstra's algorithm
reports inconsistent shortest paths after a software update that
replaced the priority queue. Argue, in 200 to 300 words, what class
of bug is most likely (priority-queue ordering, edge-weight
encoding, vertex-finalisation invariant) and what diagnostic test
would distinguish them.
\end{exercise}

\paragraph{Solution sketch.}
The change that triggered the regression is the priority-queue
replacement, so the priority-queue ordering is the prime suspect.
A common defect: the new queue's comparator orders on a different
key (insertion sequence, hash, ascending vs descending) and so the
``smallest unvisited'' vertex is no longer the one with smallest
$d(s, \cdot)$. The algorithm's correctness rests on that
invariant; violating it produces wrong distances on instances
where multiple vertices have ties or where the relaxation order
matters. Edge-weight encoding bugs (signed-vs-unsigned overflow,
units mismatch) would have manifested before the update, so they
are less likely; the finalisation-invariant bug is the same class
as the priority-queue bug (since the invariant is precisely what
the queue order enforces).

The decisive diagnostic test is a parallel run on a small,
deterministic graph for which the correct distances are known
(the worked example of figure
\ref{fig:vol02:ch17:dijkstra-trace}, for instance). Compare the
new-queue and old-queue outputs against the known answer. If the
new-queue answer differs from the known answer, the queue
ordering is at fault. A secondary test is to log the
finalisation order: with a correct queue, the finalisation order
is non-decreasing in distance. If the log shows a non-monotone
sequence, the queue is returning the wrong vertex; if the
sequence is monotone but distances are still wrong, the bug is in
the relaxation code.

\subsection*{Judgement}\pathtag{standard}

\begin{exercise}
A working engineer says: ``Asymptotic analysis is academic. The
constants matter more.'' Write a one-paragraph reply that takes the
position seriously and disagrees where disagreement is honest.
State at least one class of engineering problem in which the
asymptotic class is the only defensible figure of merit.
\end{exercise}

\paragraph{Discussion.}
The engineer's claim is half right: at fixed scale, constants
dominate asymptotic class, and a benchmark on representative
inputs trumps a theoretical bound. The discipline of asymptotic
analysis exists not because it predicts the next run, but because
it predicts how the next run scales when inputs grow. The class
of problems where the asymptotic figure is the only defensible
metric is exactly the class where the scaling matters more than the
current data: long-lived infrastructure that will see input sizes
two or three orders of magnitude beyond the present, library code
that will be called on inputs the author cannot foresee, and any
algorithm whose worst-case input determines a service-level
guarantee. In those settings, an $O(n^{2})$ algorithm with a tiny
constant is the wrong choice for a $10^{6}$-element future input
because the asymptotic class will dominate by then; the engineer
who optimised for the present has built a system that fails at
scale.

\begin{exercise}
A senior engineer dismisses the use of SAT solvers in favour of
hand-written branch-and-bound code, arguing that the SAT solver is
``a black box.'' Argue, in 200 to 300 words, when the senior
engineer's position is defensible and when it is not. Identify the
class of problem in which a modern SAT solver is the engineer's
best working tool.
\end{exercise}

\paragraph{Discussion.}
The position is defensible in three settings. First, when the
problem has known structure (a one-dimensional ordering, a metric
shortest-path, an MST) for which a polynomial-time algorithm
exists; encoding it as SAT discards the structure. Second, when
the deployment environment cannot tolerate the SAT solver's
dependencies, memory footprint, or licence; embedded systems and
hard real-time controllers are the standard examples. Third, when
the team has documented domain knowledge that a custom solver
will outperform a generic one by a wide margin (the SAT
competition winners change every year, and a domain-specific
solver tuned by experienced hands can beat the current champion
on its own instances).

The position is indefensible on combinatorial problems with no
exploited structure, where the alternative is an exponential-time
hand-written branch-and-bound that the team is reinventing badly.
Hardware-circuit verification, software bounded model checking,
combinatorial scheduling with side constraints, package-dependency
resolution, and high-level synthesis routinely pose instances with
millions of variables; modern CDCL solvers (MiniSat, Glucose,
CaDiCaL, Kissat) solve them in seconds. The engineer who
re-implements the search from scratch reinvents conflict-driven
clause learning, restarts, phase saving, variable activity
heuristics, and watched literals, all of which the SAT community
has refined for two decades. The defensible position there is to
treat the SAT solver as a thoroughly-tested algorithm and to
encode the problem clearly enough that the solver's failure mode
becomes the encoding's structure problem, not a black-box magic
problem.

\begin{exercise}
The chapter argues that combinatorial explosion is the same
archetype as the curse of dimensionality. Argue, in 200 to 300
words, whether the unification is illuminating or whether the two
phenomena are different enough that the unification obscures
useful distinctions. Take a position. Defend it.
\end{exercise}

\paragraph{Discussion.}
The unification is illuminating on the scaling axis and obscuring
on the diagnostic axis. Both phenomena share the engineer's
working observation that the size of the feasible set grows
exponentially with a parameter (variable count for combinatorial,
dimension for continuous), and both share the same three
responses (reformulate to shrink the set, exploit structure within
an exponential algorithm, accept approximation). Calling them the
same archetype is a useful pedagogy: the reader who sees one form
recognises the other immediately, and the response taxonomy
transfers.

The unification obscures the differences that matter at the
implementation level. Combinatorial spaces are unstructured by
default; continuous spaces are usually equipped with a metric and
support local search, gradient methods, and convex relaxations.
The escape from the curse of dimensionality in continuous settings
runs through Bayesian optimisation, low-rank approximation,
neural-network surrogates, and dimension reduction techniques that
have no direct combinatorial analogue. The escape from
combinatorial explosion runs through SAT/MIP/CP solvers, branch-
and-bound with combinatorial cuts, learning-augmented heuristics,
and randomisation that has no direct continuous analogue.

Position: name the archetype and use it as the scaling-literacy
hook, but resist the temptation to treat the response toolboxes as
interchangeable. The engineer who knows one toolbox should not
assume the other works the same way; the underlying mathematics is
unified, the engineering practice is not.

\begin{exercise}
A junior engineer proposes to brute-force a generative-design
search over $10^{15}$ candidate geometries by ``running it
overnight on a cluster.'' Write the response a working senior
engineer would give. State the calculation, the conclusion, and
the three responses (reformulate, exact-with-structure, approximate)
the junior engineer should consider.
\end{exercise}

\paragraph{Discussion.}
Calculation: a 1000-node cluster at $10^{10}$ ops/s per node
sustains $10^{13}$ ops/s total. Overnight is roughly $3 \times
10^{4}$ s. Total operations available: $3 \times 10^{17}$. The
search budget is $10^{15}$ candidates, each of which costs at
minimum one evaluation. If the per-candidate evaluation is a
single floating-point operation, the search fits, with a factor
of $300$ margin. If the per-candidate evaluation is anything
heavier (a finite-element solve, a CFD simulation, a structural
check), the per-candidate cost rises by $4$ to $8$ orders of
magnitude and the search falls outside the budget by the same
factor. The senior engineer's first question: \textit{what is the
per-candidate cost?}

Assume the cost is non-trivial (the usual case for generative
design). The three responses to consider:

\textit{Reformulate}: identify the symmetries and decompositions
of the design space. Geometric symmetries (parametric repetition,
reflection, translation) and constraint-driven projections often
shrink the effective space by orders of magnitude. The
$10^{15}$ may collapse to $10^{10}$ or less after the structural
quotient.

\textit{Exact with structure}: integer programming with branch and
bound, constraint programming with propagation, or SAT modulo
theories with custom propagators. The solver communities have
built engines that exploit problem structure beyond what hand-written
search reaches.

\textit{Approximate}: surrogate-model search (a neural network
trained on a sample of evaluations, used as a cheap stand-in for
the expensive per-candidate cost), Bayesian optimisation, genetic
algorithms, or modern reinforcement-learning policies. Each
accepts a probability of missing the optimum in exchange for
tractable runtime.

The conclusion the senior engineer hands back: the proposal as
stated is unfounded; pick one of the three responses and justify
it.

